\chapter{The royal cemeteries: El-Kurru and Nuri}


This chapter builds upon the history of Nubian-Egyptian relations outlined in Chapter 3 by turning to the royal necropolis of El-Kurru and Nuri. Through a quantitative and qualitative analysis of these sites, it traces the development of royal burials across the Napatan period from three distinct points: tomb architecture, object typology (excluding amulets), and amulet motifs and materials. This multi-faceted approach is structured chronologically, following the phases established previously (pre-25th Dynasty, 25th Dynasty, Early Napatan, Middle Napatan and Late Napatan) in order to correlate changes in funerary practice with the shifting relationship between Kush and Egypt. The initial analysis of tomb architecture and object types aims to establish a baseline understanding of Napatan royal burial practices and their development throughout the period. This contextual foundation is important for the subsequent and detailed study of the amulets, which seeks to identify patterns in the consumption of motifs and materials across time. Ultimately, this in-depth analysis reveals how the distinct interactions between Nubia and Egypt were materially expressed in royal amulets. These findings will provide the necessary framework for Chapter 6's comparison between royal and non-royal amulet assemblages, examining distinctions across cultural, social, chronological and regional dimensions.


\section{El-Kurru and Nuri}


The sites of El-Kurru and Nuri are well known royal sites of the Napatan period located in Upper Nubia near the 4th cataract of the Nile, with El-Kurru on the west bank and Nuri on the east bank. The sites were part of the Kushite kingdom's political center and were first excavated and recorded by Reisner from 1916 to 1920. The reports were subsequently published by Dunham in 1950 and 1955, respectively \parencites{emberlingNewFindsTwentyFifth2021, emberlingJebelBarkal201820232024}.


The El-Kurru cemetery comprises 36 tombs, including the ancestral ones, which were all plundered. According to the report, the graves were assigned an owner based on inscribed material \parencites[p.~5-12]{dunhamElKurru1950}. When no inscribed material was found, different attributions were made; from completely unknown to a possible queen of a known king. It is not clear which evidence was used to attribute titles to certain tombs. It seems that the presence of shabtis without a beard and/or its proximity to a king's pyramid ascribed the tomb to a queen. \footnote{For example, Ku. 54 and 55 do not contain shabtis and are assigned "Female, name not recovered" and "Name not recovered. Sex presumed female", respectively. The difference comes from the presence of a presumed female skeleton in Ku. 54, while Ku. 55 had no human remains. On the other hand, Ku. 71 contains at least 36 uninscribed shabtis and was assigned to "A queen (?)".} Osteological analysis of the human remains, published in 1999, revealed discrepancies between the tomb's attributed owner and the skeleton. These inconsistencies were interpreted as evidence of secondary or intrusive burials \parencites[p.~79-80]{kendallOriginNapatanState1999, beckDemographicDataHuman1999}.


The site consists of a main area, where most of the burials are located. Smaller areas northeast was tied to queens by Reisner. These include a cluster of smaller tombs (Ku. 51 to 55), while farther northeast lie the larger tombs Ku. 61 and 62. Beyond these, Ku. 71 and 72 are positioned even more remotely \parencites{emberlingNewFindsTwentyFifth2021}. Additionally, the horse burials (Ku. 200-224) are located south of Ku. 51–55. Near them, were two circular graves (Ku. 225-226), one of which contained the remains of a dog \parencites[p.~110]{dunhamElKurru1950}. Geologically, El-Kurru is located at the center of an uplifted area of land formed by ancient hard rocks. On top of this core, there are layers of much younger, flat-lying sedimentary rocks. This aspect of the site is important for the construction of the tombs, which took advantage of the geological aspects of the site for substructures to be cut into the bedrock \parencites[p.~35-38]{stearnsGeologyHydrology2016}.


New excavations at El-Kurru by the International Kurru Archaeological Project (IKAP), starting in 2013, have provided further data about the site \parencites{emberlingNewExcavationsElKurru2013}. One of the goals of the investigation was to find five structures mentioned, but not published, by Reisner in his field diary and re-discovered by Timothy Kendall. These are two funerary temples (Ku. 1100 and 1500), a rock-cut well (Ku. 1400), a city wall and gate (Ku. 1300) and a portion of a fortification wall (Ku. 1200). The excavations located four structures, with the exception of the fortification wall \parencites[p.~47-49]{kendallOriginNapatanState1999}[p.~65-67]{ emberlingNewFindsTwentyFifth2021}.


Excavations in the larger mortuary temple (Ku. 1500) uncovered columns similar to those in Ptolemaic Alexandria, seven underground rock-cut chambers---four with secure doors---and two outer rooms. Based on the lack of archaeological objects, unfinished blocks in some rooms and concern for security, the excavators first suggested that the structure could have been designed as a protected storeroom for valuable objects \parencites{chengFivesidedCorinthianCapitals2013}[p.~65-67]{emberlingExcavationPyramidKu2015}{karbergRemarksStonemasonsMarks2015}. However, in a later publication about the smaller temple, it was alluded that both structures were used first as quarries and then expanded with the rock-cut rooms, but were “not finished and used for its original intended purpose” \parencites[p.~38-39]{anstisLateNapatanRockcut2016}. Both of these structures were not related to the main occupation of the site, from around 750 to 664 BCE, but were a build during the 4th century BCE. The other identified structures---the city wall, well and gate (Ku. 1300)---actually date to the Medieval Christian period \parencites[p.~43-44]{emberlingNewExcavationsElKurru2013}{skuldbølMedievalFortificationSettlement2016}.


The recent excavations have pointed that after the initial occupation of the burial site, there was an effort to return to it with both temples and the unfinished pyramid Ku. 1. One suggestion is that competition between ruling families in Napata during the 4th century BCE caused the transfer of the royal burial site to nearby Meroe \parencites{emberlingExcavationPyramidKu2015}. Ultimately, at the present time, there are no traces of a settlement in El-Kurru, contrary to what Reisner had first suggested \parencites[p.~68]{emberlingNewFindsTwentyFifth2021}.


At Nuri, there are 73 tombs in total and the owner assignment seems to have followed the same logic as in El-Kurru. All tombs were also plundered, some even when the cemetery was still in use \parencites[p.~4]{dunhamNuri1955}. The first king to move his burial place from El-Kurru to Nuri was Taharqa (Nu. 1). Scholars have posited different theories as to why Taharqa decided to build his pyramid in a different location: due to astronomical characteristics of the site or to the impossibility of being interred in Egypt at the time. Afterwards, Nuri became the new royal burial site for succeeding kings (with the exception of Tanwetamani, his queen and the two unoccupied later pyramids) \parencites{kendallWhyDidTaharqa2008, lohwasserDoublingDoubleKingdom2019}.


Some aspects of the tombs, unfortunately, are not totally known given the state of preservation at the time of the excavations. For instance, some superstructures from both sites were considered unknown, since the reports state that no trace of them was found. However, it is important to consider that could be due to their removal or destruction and not necessarily that there was no superstructure at all \parencites{mohammed–aliUniversityDongolaArchaeological2016}.\footnote{The Nuri report even notes that the superstructures of Nu. 53 and 45 were destroyed by later Christian constructions \parencites[p.~53, 149]{dunhamNuri1955}.}


Regarding El-Kurru, Reisner and Dunham analyzed the superstructures as an evolution in a Egyptocentric way: starting with the tumulus, then the mastaba and, finally, the pyramid \parencites[p.~123-124]{dunhamElKurru1950}. Reisner divided the tombs according to the idea that a mastaba would have a burial pit and a pyramid would have underground chambers. However, by analyzing the ground plans of the mastabas, Lacovara \textcite[p.~146]{lacovaraTumulusPyramidDevelopment2018} has suggested that they were actually early pyramids and there was no mastaba stage at El-Kurru. This idea is also substantiated by the lack of mastabas in Nubia during the New Kingdom \parencites[p.~27]{kendallOriginNapatanState1999}. According to \textcites[p.~11, 26-27]{kendallOriginNapatanState1999}, Reisner had considered this possibility, but, likely due to his belief that the Nubian tombs had to have the same evolution as the Egyptian ones, decided that the square El-Kurru tombs must have been mastabas. Given the evidence, it seems likely that, at El-Kurru, there were two types of superstructure: the earlier tumulus and the later pyramid. Similarly, Nuri has a superstructure described in the report as a “stepped pyramid or mastaba” (Nu. 37). I would argue that this structure should also be classified as a pyramid, applying the same criteria used for the mastabas at El-Kurru.


The tombs at both sites represent different phases of the Napatan period and will be described separately in the following sections in order to examine particular chronological variations. Additionally, each site and phase includes tombs from varying hierarchies within the royal sphere, ranging from ruling kings to prominent queens and minor wives, following Reiner's designation. This hierarchical distinction will also be examined within each phase.


\section{Pre-25th Dynasty}


\subsection{Tomb structure}


The thirteen burials from the pre-25th Dynasty phase are found exclusively at El-Kurru. Four tombs were associated with either kings/queens, six to kings, two to queens and one to a minor wife. Their superstructures consist of tumuli, pyramids and one unknown structure---likely destroyed (Figure \ref{fig:pre-royal-tombs}). The tumuli are solely associated with kings or queens, depending on whether the long or short chronology is applied. Pyramids are linked with both kings and queens, while the unidentified superstructure belonged to a minor wife. All of the tombs included a pit as the sole substructure.


Other architectural elements include chapels, a side chamber and holes for a burial bed (Table \ref{tab:pre-tomb-elements}). It is likely that some of these features did not survive. Six out of the thirteen tombs included a funerary chapel, which were mostly found in later graves. The chapels could be a reintroduction in Nubian burials, as they were attested in Kerma tumuli and elite C-Group burials \parencites[p.~2221-223]{bonnetEasternCemeteryKerma2020}[p.~162]{hafsaasCGroupPeopleLower2020}. In fact, one tumulus (Ku. Tum. 6) included an attached chapel. Nevertheless, most chapels were associated with pyramids. Nine of the tombs included a side chamber, where the burial took place. This feature was present in both tumuli and pyramids from kings and queens. Three of the side chambers included four holes on the floor for the placement of a burial bed. Although lacking these cutouts, Ku. 13 yielded decayed wooden fragments of a bed, indicating that burial beds could be included without the holes, or that the holes did not survive. Out of the three tombs with bed holes, only one is a tumulus, while the other two were in pyramids. Additionally, this is the only phase of the Napatan period in which not only  coffin benches were absent, but also coffin fittings or eye inlays. \parencites[p.~51]{dunhamElKurru1950}.


\begin{figure}[H]
	\centering
	\includegraphics[width=1\textwidth]{pre-royal-tombs}
	\caption{Distribution of superstructure and substructure types by occupant's title in pre-25th dynasty royal burials.}
	\label{fig:pre-royal-tombs}
\end{figure}



{\small
	\begin{xltabular}{\textwidth}{X X X X X X}
		\caption{Distribution of architectural elements by tomb owner's title in pre-25th dynasty royal burials.}
		\label{tab:pre-tomb-elements} \\
		\toprule
		\textbf{Tomb} & \textbf{Super- and substructure} & \textbf{Owner} & \textbf{Chapel} & \textbf{Side chamber} & \textbf{Bed holes and/or coffin bench} \\
		\midrule
		\endfirsthead
		\multicolumn{6}{c}{\textit{Continued from previous page}} \\
		\toprule
		\textbf{Tomb} & \textbf{Super- and substructure} & \textbf{Owner} & \textbf{Chapel} & \textbf{Side chamber} & \textbf{Bed holes and/or coffin bench} \\
		\midrule
		\endhead
		\bottomrule
		\endlastfoot
		Ku. Tum. 1 & tumulus/pit & king & -- & yes & -- \\
		\cmidrule(r){1-6}
		Ku. Tum. 2 & tumulus/pit & king & -- & -- & -- \\
		\cmidrule(r){1-6}
		Ku. Tum. 4 & tumulus/pit & king or queen & -- & -- & -- \\
		\cmidrule(r){1-6}
		Ku. Tum. 5 & tumulus/pit & king or queen & -- & yes & -- \\
		\cmidrule(r){1-6}
		Ku. Tum. 6 & tumulus/pit & king & yes & yes & -- \\
		\cmidrule(r){1-6}
		Ku. 19 & tumulus/pit & king or queen & -- & yes & bed holes \\
		\cmidrule(r){1-6}
		Ku. 14 & pyramid/pit & king & yes & yes & bed holes \\
		\cmidrule(r){1-6}
		Ku. 13 & pyramid/pit & king or queen & yes & yes & -- \\
		\cmidrule(r){1-6}
		Ku. 11 & pyramid/pit & king & yes & yes & -- \\
		\cmidrule(r){1-6}
		Ku. 10 & pyramid/pit & queen & yes & yes & bed holes \\
		\cmidrule(r){1-6}
		Ku. 9 & pyramid/pit & king & yes & yes & -- \\
		\cmidrule(r){1-6}
		Ku. 23 & pyramid/pit & queen & -- & -- & -- \\
		\cmidrule(r){1-6}
		Ku. 21 & unknown/pit & minor wife & -- & -- & -- \\
	\end{xltabular}
}


\subsection{Object types}


\begin{figure}[H]
	\centering
	\includegraphics[width=1\textwidth]{pre-royal-objs}
	\caption{Distribution of object types in pre-25th dynasty royal tombs by occupant's title.}
	\label{fig:pre-royal-objs}
\end{figure}


In pre-25th Dynasty tombs, beads are the most common object, with 532 items in king's and 103 in king/queen's tombs (Figure \ref{fig:pre-royal-objs}). Weapons and ceramic and non-ceramic vessels were also popular, present in the tombs of both kings and queens. Other object types were much rarer. Only one offering table and one scarab were found, both in a king/queen tomb. Shabtis were present in all types of burials: one with a king, two with a king/queen and one with a queen. King/queen's tombs were the most diverse, containing twelve different object types. Kings' burials included eight, whereas queens' contained only the sole shabti.


\subsection{Amulets}


\begin{figure}[H]
	\centering
	\includegraphics[width=1\textwidth]{pre-royal-motifs}
	\caption{Distribution of amulet motifs in pre-25th dynasty royal tombs.}
	\label{fig:pre-royal-motifs}
\end{figure}


\begin{wrapfigure}{l}{0.4\textwidth}
	\centering
	\includegraphics[width=0.50\linewidth, height=4cm, keepaspectratio]{pre-amulet-1}
	\includegraphics[width=0.50\linewidth, height=3.8cm, keepaspectratio]{pre-amulet-2}
	\caption{Amulet of an inscribed gold nugget and double-sided amulet of a Pataikos and hawk-headed deity (\cite[pp.~16, fig.~2d, pl.~LIIB]{dunhamElKurru1950})}
	\label{fig:pre-amulets}
\end{wrapfigure}


The sole pre-25th Dynasty tomb containing amulets is Ku. Tum. 2, associated with a king. The tomb contained 29 amulets, all made of gold (Figure \ref{fig:pre-royal-motifs}). Twenty-seven of these amulets depicted the \textit{udjat} eye and were all strung together into a necklace, between garnet beads. The \textit{udjat}, a very common amulet, was used in Nubia since the New Kingdom, being found in burials of different sites \parencites[p.~162]{lemosMaterialCultureColonization2020}. Another amulet was a natural gold nugget inscribed with hieroglyphics that read "Amun, lord of life, that which he gives is beautiful, Amun, lord…" \parencites[p.~25]{kendallOriginNapatanState1999} (Figure \ref{fig:pre-amulets}, left). This evidence suggests not only the presence of people who knew the Egyptian language in Nubia during this period, but also the integration of the cult of Amun into the royal sphere. The last amulet depicts a Pataikos figure with a hawk-headed deity on the reverse (Figure \ref{fig:pre-amulets}, right). As we will see later, double-sided deity amulets and of hawk-headed figures are also found in other phases of the Napatan period. Thus, the presence of one of these amulets in one of the ancestral tombs may indicate its possible origin. Both the nugget and double-sided amulet were strung together in a necklace with gold beads. Typologically, most amulets depict symbols, while one represented deities. The absence of amulets made of other materials is interesting as gold was a local commodity in Nubia, indicating the possibility that acquisition of imported materials for amulets was limited during this period \parencites{klemmGoldGoldMining2013}.


\section{25th Dynasty}


\subsection{Tomb structure}


\begin{figure}[H]
	\centering
	\includegraphics[width=1\textwidth]{25-royal-tombs}
	\caption{Distribution of superstructure and substructure types by occupant's title in 25th dynasty royal burials.}
	\label{fig:25-royal-tombs}
\end{figure}


25th Dynasty tombs are found in both El-Kurru and Nuri (Figure \ref{fig:25-royal-tombs}). In total, there are 27 tombs, 20 at El-Kurru and seven at Nuri. Of those, sixteen belong to queens, six to kings and five to minor wives. At El-Kurru, tomb superstructures include pyramids and unidentified ones, which could be paired with chambers or a burial pit. The most common structure are pyramids with chambers (8), followed by the unknown superstructures---most likely destroyed pyramids---with pits (6). Conversely, four tombs feature an unidentified superstructure with chambers, while there are two pyramids with pits. Similar to the pre-25th Dynasty tombs, minor wives' tombs at El-Kurru also feature an unknown superstructure with a burial pit. Tombs of queens have more diverse structures, such as pyramids with chambers (4) or a pit (1), unidentified superstructures with chambers (4) or a pit (4). On the other hand, kings' tombs have only pyramids as superstructures featuring either chambers (4) or pits (1). Even if we consider all of the unknown superstructures as possible destroyed pyramids, queens and kings' tombs would still feature a pyramid with both substructures, while the minor wives' burials would include solely a burial pit.


The unknown superstructures at Nuri were described as “destroyed” in Dunham's report, with some leaving faint traces, while most had vanished entirely. However, given the predominance of pyramids at the site, it is likely that these structures were also destroyed pyramids, like in El-Kurru.	The later burials at Nuri reveal a significant change in the funerary sphere. Pits are no longer used as a substructure---all tombs only include chambers. However, the superstructure of the minor wives' tombs were still unidentified and queen's were either an unknown superstructure or a pyramid. 

25th Dynasty tomb superstructures at Nuri include an innovation: the double pyramid. This structure is characterized by two pyramids, one inside the other. Only two tombs were made in this fashion during this phase: Taharqa's tomb, Nu. 1, the largest in the cemetery, and  his queen Atakhebasken (Nu. 36). Taharqa's structure was built in three different stages, with the outer pyramid constructed after his burial \parencites[p.~118-119]{kendallWhyDidTaharqa2008}. Queen Atakhebasken's pyramid appears to have been enlarged at a later date as well \parencites[p.~108]{welsbyKingdomKushNapatan1996}. Therefore, both were initially designed as conventional pyramids, with Nu. 36 potentially expanded due to the queen's association with Taharqa. Applying the hypothesis that all unidentified superstructures were originally pyramids, and acknowledging that even the double pyramids began as single pyramids, we can posit a consistent funerary pattern at Nuri: all burials, irrespective of the occupant's title, feature pyramids with underground chambers.


The tomb architectural features changed throughout the 25th Dynasty (Table \ref{tab:25-tomb-elements}). hapels were not consistently incorporated into royal tombs during this period, appearing in some kings' and queens' burials but being entirely absent from those of minor wives, without a clear pattern. However, this may be due to the destruction of these structures.


Among the earliest tombs at El-Kurru (Ku. 8 to 53), only Piankhy's (Ku. 17) had a stairway entrance leading to a single chamber. In contrast, Tabiry's tomb (Ku. 53), his queen, featured a particular design that could be interpreted as a type of "shaft" entrance with a single step into the burial chamber (Figure \ref{fig:25-tabiry-tomb}). This distinctive element has no known counterparts in royal burials and raises the possibility that the substructure was originally intended as a pit with a side chamber---similar to those from the ancestral tumuli---, suggesting that the "shaft entrance" may not, in fact, be an entrance at all. Whether the intention was to produce a tomb similar to the ancestral ones or it was just a coincidence driven by other factors, such as geographical limitations or personal preference, is unknown.


Queen Tabiry's tomb was located in an area of the El-Kurru cemetery near other tombs: Ku. 51, 52, 54 and 55. All of the other tombs substructures were pits without side chambers and belonged to other queens and minor wives of Piankhy. Interestingly, stairways at El-Kurru are believed to have been introduced during the reign of Piankhy as it allowed the superstructure to be built over the chambers while the owner was still alive \parencites[p.~623]{lohwasserNapatanNecropoleisBurial2019}. However, the majority of the burials of Piankhy's queens and minor wives did not contain this type of entrance. The exception is Ku. 4, the tomb of queen Khensa, whose tomb featured a stairway and two chambers, more than Piankhy's single-chambered tomb. Shabataqo's tomb (Ku. 18) established stairs as the definitive type of entrance in royal burials.

Among the burials with chambers as a substructure, the number appear to follow a pattern. The first tomb with chambers, Piankhy's Ku. 17, features only one chamber. From Shabataqo's burial, all kings' tombs include two chambers, whereas some queens' tombs could have one or two chambers. All the tombs of the minor wives featured only one chamber. Decoration either on chapel walls or the chambers was inconsistent. Out of the total seven decorated graves, four were of queens and three of kings. As with the chapels, this may be due to preservation issues.


The majority of early tombs at El-Kurru displayed both a coffin bench and cutouts for a burial bed, a tradition that ceases with the burial of a queen of Taharqa (Nu. 35), where only coffin benches were included. In Taharqa's tomb, Dunham described a rectangular sinkage on the middle of the burial chamber, presumed for a wooden coffin, since no trace was found \parencites[p.~9]{dunhamNuri1955}. In fact, Taharqa's tomb exhibits an unique substructure, emulating the Osireion of Seti I in Abydos and identifying the king with the god Osiris (Figure \ref{fig:25-taharqa-tomb}) \parencites[p.~120]{kendallWhyDidTaharqa2008}[p.~76-77]{lohwasserDoublingDoubleKingdom2019}.


Direct evidence for coffins is scarce, with some burials providing hints of its use. In Piankhy's queen Neferukakashta (Ku. 52) bench, Reisner noted remains of wood and plaster bearing hieroglyphs. Further evidence stems from decorative elements, such as eye inlays and coffin fittings---found in the tombs of queen Khensa (Ku. 3), queen Naparaye (Ku. 4), Shabaqo (Ku. 15) and three of Piankhy's queens (Ku. 51, 52, 53), all of which included a coffin bench and bed holes.


As for burial beds, four tombs contained either bed legs or containers for them. Notably, the burial of a queen of Shabataqo (Ku. 72) featured two decorated bronze bed legs \parencites[p.~104]{dunhamElKurru1950}.. Since only the legs survived, it is possible that the rest of the bed was made of wood. At queen Naparaye's and another minor wife of Taharqa's tombs (Ku. 4 and Nu. 35), Reisner found alabaster boxes for bed legs \parencites[p.~31]{dunhamElKurru1950}[p.~19]{dunhamNuri1955}. Ku. 4 featured a coffin bench and bed cutouts, whereas Nu. 35 included only a coffin bench. This suggests that burial beds may have been used independently of the cutouts, or that the holes did not survive. Consequently, later tombs without cutouts could potentially have contained wooden beds placed on the coffin bench.


{\small
	\begin{xltabular}{\textwidth}{X X X X X X X}
		\caption{Distribution of architectural elements by tomb owner's title in 25th dynasty royal burials.}
		\label{tab:25-tomb-elements} \\
		\toprule
		\textbf{Tomb} & \textbf{Owner} & \textbf{Chapel} & \textbf{Type of entrance} & \textbf{N. of chambers} & \textbf{Decoration} & \textbf{Bed holes and/or coffin bench} \\
		\midrule
		\endfirsthead
		\multicolumn{7}{c}{\textit{Continued from previous page}} \\
		\toprule
		\textbf{Tomb} & \textbf{Owner} & \textbf{Chapel} & \textbf{Type of entrance} & \textbf{N. of chambers} & \textbf{Decoration} & \textbf{Bed holes and/or coffin bench} \\
		\midrule
		\endhead
		\bottomrule
		\endlastfoot
		Ku. 8 & king & yes & -- & -- & -- & -- \\
		\cmidrule(r){1-7}
		Ku. 7 & queen & yes & -- & -- & -- & -- \\
		\cmidrule(r){1-7}
		Ku. 20 & minor wife & -- & -- & -- & -- & -- \\
		\cmidrule(r){1-7}
		Ku. 17 & king & yes & stairs & 1 & -- & both \\
		\cmidrule(r){1-7}
		Ku. 22 & queen & -- & -- & -- & -- & -- \\
		\cmidrule(r){1-7}
		Ku. 55 & queen & -- & -- & -- & -- & both \\
		\cmidrule(r){1-7}
		Ku. 54 & minor wife & -- & -- & -- & -- & both \\
		\cmidrule(r){1-7}
		Ku. 52 & queen & -- & -- & -- & -- & both \\
		\cmidrule(r){1-7}
		Ku. 51 & queen & -- & -- & -- & -- & both \\
		\cmidrule(r){1-7}
		Ku. 53 & queen & -- & shaft? & 2? & -- & both \\
		\cmidrule(r){1-7}
		Ku. 4 & queen & yes & stairs & 2 & -- & both \\
		\cmidrule(r){1-7}
		Ku. 18 & king & yes & stairs & 2 & -- & both \\
		\cmidrule(r){1-7}
		Ku. 72 & queen & -- & stairs & 1 & -- & both \\
		\cmidrule(r){1-7}
		Ku. 15 & king & yes & stairs & 2 & yes & both \\
		\cmidrule(r){1-7}
		Ku. 71 & queen & -- & stairs & 1 & -- & both \\
		\cmidrule(r){1-7}
		Ku. 62 & queen & -- & stairs & 1 & yes & both \\
		\cmidrule(r){1-7}
		Nu. 1 & king & -- & stairs & 2 & yes & coffin sinkage? \\
		\cmidrule(r){1-7}
		Ku. 3 & queen & -- & stairs & 2 & -- & both \\
		\cmidrule(r){1-7}
		Nu. 35 & queen & yes & stairs & 2 & -- & coffin bench \\
		\cmidrule(r){1-7}
		Nu. 36 & queen & yes & stairs & 2 & yes & -- \\
		\cmidrule(r){1-7}
		Ku. 16 & king & yes & stairs & 2 & yes & -- \\
		\cmidrule(r){1-7}
		Ku. 5 & queen & yes & stairs & 2 & yes & coffin bench \\
		\cmidrule(r){1-7}
		Ku. 6 & queen & -- & stairs & 2 & yes & coffin bench \\
		\cmidrule(r){1-7}
		Nu. 59 & queen & -- & stairs & 1 & -- & coffin bench \\
		\cmidrule(r){1-7}
		Nu. 74 & minor wife & -- & stairs & 1 & -- & coffin bench \\
		\cmidrule(r){1-7}
		Nu. 77 & minor wife & -- & stairs & 1 & -- & -- \\
		\cmidrule(r){1-7}
		Nu. 80 & minor wife & -- & stairs & 1 & -- & -- \\
	\end{xltabular}
}


\begin{figure}[H]
	\centering
	\includegraphics[width=0.6\textwidth]{25-tabiry-tomb}
	\caption{Tomb Ku. 53, attributed to queen Tabiry, main wife of Taharqa \parencites[p.~86, fig.~29a]{dunhamElKurru1950}}
	\label{fig:25-tabiry-tomb}
\end{figure}


\begin{figure}[H]
	\centering
	\includegraphics[width=0.6\textwidth]{25-taharqa-tomb}
	\caption{The substructure of Taharqa's tomb, Nu. 1, showing the sinkage on the middle of the burial chamber  \parencites[p.~8, fig.~2b]{dunhamNuri1955}}
	\label{fig:25-taharqa-tomb}
\end{figure}


\subsection{Object types}


The assemblage of object types from the 25th Dynasty is significantly more diverse than the previous phase (Figures \ref{fig:25-royal-objs-el-kurru}, \ref{fig:25-royal-objs-nuri}). It includes 13 new categories: bead nets, burial containers, canopics, heart scarabs, mummy trappings, games and toys, pebbles, shrines, staffs, stelae, tools and toiletries. Some of these object types will continue to appear in later phases, even as pebbles and games and toys were only found in this phase.


The most common object types at both El-Kurru and Nuri are shabtis and individual beads. The distribution of these items varied by site and owner. At El-Kurru, shabtis were found in queens' (1214) and kings' tombs (939). At Nuri, the distribution was different: shabtis were mostly found in the tombs of kings (1070), followed by those of minor wives (520) and, then, queens (238). While individual beads are the second most popular object at both sites--- 216 in minor wives' tombs, 171 in queens', and 80 in kings' at El-Kurru---the distribution in Nuri was different. Beads were most prevalent in queens' tombs (100), followed by kings' (69), and were totally absent from the graves of minor wives'. However, in both sites, adornments were featured in all types of burials.


\begin{figure}[H]
	\centering
	\includegraphics[width=\textwidth, height=0.6\textheight, keepaspectratio]{25-royal-objs-el-kurru}
	\caption{Distribution of object types in 25th dynasty royal tombs from El-Kurru by occupant's title.}
	\label{fig:25-royal-objs-el-kurru}
\end{figure}


Ceramic and non-ceramic vessels were popular objects as well. Included in the non-ceramic vessels category are round bronze trays only attested at El-Kurru during this phase, without parallels in Egypt. Analysis of a tray from Queen Khensa (Ku. 4), currently at the Museum of Fine Arts, Boston, revealed elaborate decoration, contact with organic materials and incomplete inscriptions referring to use of natron in purification \parencites[p.~499-500]{gänsickeExaminationThreeLargeScale2015}. Unlike the tombs from the pre-25th Dynasty phase, no weapons were found at El-Kurru, but they were present in kings' tombs at Nuri, albeit in small quantities. 


\begin{figure}[H]
	\centering
	\includegraphics[width=\textwidth, height=0.6\textheight, keepaspectratio]{25-royal-objs-nuri}
	\caption{Distribution of object types in 25th dynasty royal tombs from Nuri by occupant's title.}
	\label{fig:25-royal-objs-nuri}
\end{figure}


Heart scarabs are an interesting case. Two instances of this object were present in just one tomb, Tanwetamani's (Ku. 16). This is paralleled by a burial in Aniba from the New Kingdom, which also included two heart scarabs, a restricted object during the period. Lemos has interpreted this case as either the owner's display of access to foreign goods, or his connection to people in life who could gift him two instances of this restricted object \parencites[p.~260]{lemosHeartScarabsOther2021}. At Nuri, only one heart scarab was found, in the tomb of Tanwetamani's queen Maloqe. Since there were only these three heart scarabs in tombs of the 25th Dynasty, it is possible that either these objects continued to be restricted in the Napatan period, even for royalty, or Tanwetamani restarted the use of heart scarabs that only lasted in his reign. Regarding the canopics category, two of the jars in El-Kurru were “dummies”, meaning that the were solidly made, with not separation between the body of the jar and lid, and, therefore, could not be used to store the organs after mummification. Both were found in the tomb of Shabataqo (Ku. 18) (Dunham, 1950, p. 68). As Taharqa's and his queens canopic jars were usable, it is assumed that during his reign, canopic jars were again used for their original purpose \parencites[p.~299]{astonBurialAssemblagesDynasty2009}.
 
 
Beyond the frequency, the diversity of object types also varied by hierarchy. Queens possessed the most diverse assemblages at both El-Kurru (25 types) and Nuri (16 types). Kings' tombs featured less object types (15 at El-Kurru and 14 at Nuri), while minor wives had the least diversity, with only 7 in El-Kurru and 8 in Nuri. This pattern could suggest that either queens were furnished with a wider diversity of items than even kings, or that their pyramids were less looted. Some object types are only present at El-Kurru---pebbles, statuettes, tools, jar stoppers, games and toys, stelae, bead nets, clothing and shrines---, which could be due to the bigger number of tombs at that site. Three object types are exclusive to Nuri: staffs, weapons and mummy trappings. Notably, all preserved mummy trappings were made of gold, silver or electrum, which likely indicates they were overlooked by the looters.
 
 
 \subsection{Amulets}
 
 
 \subsubsection{Amulet distribution}
 
 
 {\small
 	\begin{xltabular}{\textwidth}{X X X X X X}
 		\caption{Distribution of amulets in 25th Dynasty royal tombs by owner's title, including total counts of amulets and tombs, tombs containing amulets and amulet frequency (\%) and median.}
 		\label{tab:25-royal-distro} \\
 		\toprule
 		\textbf{Owner} & \textbf{Total number of tombs} & \textbf{Total number of amulets} & \textbf{N. of tombs with amulets} & \textbf{\% of tombs with amulets} & \textbf{Median n. of amulets} \\
 		\midrule
 		\endfirsthead
 		\multicolumn{6}{c}{\textit{Continued from previous page}} \\
 		\toprule
 		\textbf{Owner} & \textbf{Total number of tombs} & \textbf{Total number of amulets} & \textbf{N. of tombs with amulets} & \textbf{\% of tombs with amulets} & \textbf{Median n. of amulets} \\
 		\midrule
 		\endhead
 		\bottomrule
 		\endlastfoot
 		King & 6 & 172 & 6 & 100\% & 10 \\
 		\cmidrule(r){1-6}
 		Queen & 16 & 285 & 11 & 70\% & 4 \\
 		\cmidrule(r){1-6}
 		Minor wife & 5 & 8 & 2 & 40\% & 0 \\
 		\midrule
 		\textbf{Total} & \textbf{27} & \textbf{465} & \textbf{19} & -- & -- \\
 	\end{xltabular}
 }
 
 
 There is a total of 465 amulets distributed across 19 tombs (Table \ref{tab:25-royal-distro}). 70\% of queens' tombs contained amulets, comprising a total of 285 (61\% of the total). While all kings' burials contained amulets, they yielded a smaller combined total of 172. Among the tombs of minor wives, only two of the five contained amulets, with 8 examples found. Although queen's were interred with the biggest absolute number of amulets, it is important to look at the median number of amulets, i.e. the middle number. In the case of amulets, the median is more informative than the average, as many tombs contained no amulets at all, skewing the data. Therefore, the median will be used in order to more accurately interpret the distribution of amulets across the different contexts. In this phase, although queen's were interred with the biggest absolute number of amulets, the median was higher with kings (10 against 4, respectively). Wives, on the other hand, had a median of zero. This means that kings were more likely to be buried with amulets.
 
 
 \subsubsection{Amulet types}
 
 
  The distribution of amulet types in 25th Dynasty tombs differs significantly between El-Kurru and Nuri, as well as the occupant's title (Figure \ref{fig:25-royal-types}). The variations are partly due to the greater number of tombs at El-Kurru, which yielded 396 amulets compared to only 69 from Nuri.
  
 
 \begin{figure}[H]
 	\centering
 	\begin{minipage}{1\textwidth}
 		\centering
 		\includegraphics[width=\textwidth, height=8cm, keepaspectratio]{25-royal-types-el-kurru}
 	\end{minipage}
 	
 	\begin{minipage}{1\textwidth}
 		\centering
 		\includegraphics[width=\textwidth, height=8cm, keepaspectratio]{25-royal-types-nuri}
 	\end{minipage}
 	\caption{Distribution of amulet types in 25th Dynasty royal tombs by occupant's title.}
 	\label{fig:25-royal-types}
 \end{figure}
 
 
 At El-Kurru, symbol amulets are the most numerous, predominantly found in kings' tombs (143), compared to 96 in queens' and only two in minor wives'. Conversely, queens' burials contained the majority of deity symbols (126)---the second most common type---far surpassing those of kings' (2) and minor wives' (4). Furthermore, animal amulets were exclusive to queens (15) and kings (2). The few amulets depicting nature elements, objects and miscellaneous motifs were found solely with queens. The distribution shows that queens possessed not only the most amulets, but also the greatest diversity, as with object types and perhaps for the same reasons.
 
 
 A similar pattern emerges at Nuri, where symbol amulets were also the most frequent. Queens' tombs, however, held the majority (29), compared to kings' (22) and minor wives' (2). Nature amulets were the second most common type, but exclusively found with queens (6). While animal amulets were found with both kings (3) and queens (2), deity amulets were attested only in female burials (three with a queen and one with a minor wife). The only object amulet was found in a queen's burial. Similar to El-Kurru, queens tombs at Nuri contained the most diverse assemblage of amulet types.
 
 
 \subsubsection{Amulet motifs}
 
 
 This section examines the specific motifs of each type in more detail. Given the high number of unique motifs, they are presented in separate graphs based on their broader typological groups for better visualization.
 
 
 \begin{figure}[H]
 	\centering
 	\includegraphics[width=1\textwidth]{25-royal-deities-el-kurru}
 	\caption{Distribution of amulets with deity motifs in 25th Dynasty royal tombs by occupant's title at El-Kurru.}
 	\label{fig:25-royal-deities-el-kurru}
 \end{figure}
 
 
 The amount of deity amulets at El-Kurru is very high, with 31 different motifs (Figure \ref{fig:25-royal-deities-el-kurru}). Burials of kings contained only two of these amulets; one depicting a winged goddess and another, a lion-headed deity. Minor wives' were buried with four deity amulets: two Pataikos, one Hathor and one Amun. The rest of the motifs were all found with queens, 30 in total. The most popular deity is, by far, the Pataikos, with 38 instances distributed across female burials.
 
 
 The deity amulets show strong Egyptian influence in their motifs, yet many were adapted to reflect the Kushite context. One of these amulets are those depicting winged deities or female figures. These amulets occur in only three burials, those of Piankhy's queens: an unnamed one (Ku. 51), Neferukakashta (Ku. 52) and Tabiry (Ku. 53). Fifteen of these amulets depict unknown anthropomorphic female figures, while five have leonine heads. A lion-headed deity also appears as an aegis at Tanwetamani's tomb (Ku. 16). The winged goddesses are distinguished in the way they are portrayed, with prominent thighs and breasts, similar to the way Meroitic queens will be depicted later \parencites[p.~65]{lohwasserRoleStatusRoyal2020}. They also carry symbols of Egyptian origin, such as the sun disk, \textit{uraeus}, and the double plume crown. All of the amulets are different, displaying diverse elements and even perspectives (Figure \ref{fig:25-goddesses}). Most of the female figures are depicted frontally, with their upper arms raised and forearms downward, creating an arch over their wings.
 
 
\begin{figure}[H]
	\centering
	\begin{minipage}{0.75\textwidth}
		\centering
		\begin{minipage}{0.3\textwidth}
			\includegraphics[width=\textwidth, height=4.5cm, keepaspectratio]{25-24.616}
		\end{minipage}
		\hspace{-8mm}%
		\begin{minipage}{0.3\textwidth}
			\includegraphics[width=\textwidth, height=4.5cm, keepaspectratio]{25-24.683}
		\end{minipage}
		\hspace{-8mm}%
		\begin{minipage}{0.3\textwidth}
			\includegraphics[width=\textwidth, height=4.5cm, keepaspectratio]{25-24.1800}
		\end{minipage}
		\vspace{0mm}%
		\begin{minipage}{0.3\textwidth}
			\includegraphics[width=\textwidth, height=4.5cm, keepaspectratio]{25-sudan-1}
		\end{minipage}
		\hspace{-10mm}%
		\begin{minipage}{0.3\textwidth}
			\includegraphics[width=\textwidth, height=4.5cm, keepaspectratio]{25-24.707}
		\end{minipage}
		\hspace{-8mm}%
		\begin{minipage}{0.3\textwidth}
			\includegraphics[width=\textwidth, height=4.5cm, keepaspectratio]{25-sudan-2}
		\end{minipage}
	\end{minipage}%
	\hspace{-22mm}
	\begin{minipage}{0.2\textwidth}
		\raggedright
		\captionof{figure}{Amulets from 25th Dynasty El-Kurru depicting anthropomorphic and lion-headed female figures (from top to bottom and left to right: MFA 24.1800, 24.683, SNM 2014, MFA 24.707, 24.616, SNM 2015).}
		\label{fig:25-goddesses}
	\end{minipage}
\end{figure}


In fact, winged goddesses, leonine and anthropomorphic, were featured in the reverse of many amulets (Figure \ref{fig:25-winged-back}). Figure \ref{fig:25-royal-2deities} illustrates the correlation of motifs in multifaceted amulets. Nine Pataikos amulets included another motif on its reverse, including symbols, apes, the god Shu and the winged goddesses---one anthropomorphic and two with a lion's head. The latter also appears on the back of amulets depicting Egyptian deities such as the triad Amun, Isis and Horus, as well as more Kushite designs such as a ram-headed dwarf and a winged Pataikos. These amulets all come from the same group of tombs belonging to Tanwetamani’s queens (Ku. 51, 52 and 53).


\begin{figure}[H]
	\centering
	\includegraphics[width=1\textwidth, height=4.5cm, keepaspectratio]{25-24.690}
	\includegraphics[width=1\textwidth, height=4.5cm, keepaspectratio]{25-24.690-2}
	\caption{Example of a Pataikos amulet featuring a lioness winged goddess on the reverse, from Ku. 53 (MFA 24.690)}
	\label{fig:25-winged-back}
\end{figure}


\begin{figure}[H]
	\centering
	\includegraphics[width=1\textwidth]{25-royal-2deities}
	\caption{Relation of 25th Dynasty multifaceted royal amulets.}
	\label{fig:25-royal-2deities}
\end{figure}


Two other amulets display winged figures: a winged Pataikos and an unknown ram-headed god. Conversely, a non-winged Pataikos displays a ram head. Both of these ram-headed amulets feature a double \textit{uraeus}, a marker of Kushite rule \parencites[p.~6-7]{törökRoyalCrownsKush1987} (Figure \ref{fig:25-royal-pataikos}). Interestingly, two other amulets from the same tombs (Neferukakashta and Tabiry's) depict unidentified ram-headed deities, but without wings. Amun was depicted with a ram head by the Kushites and was considered the local form of the god \parencites[p.~344]{gaboldeAmunCultIts2020}. It is possible that this amulets were intended to establish an association between Amun and other gods, such as the Pataikos, depicted as a dwarf.


\begin{figure}[H]
	\centering
	\begin{minipage}{0.3\textwidth}
		\centering
		\includegraphics[width=\textwidth, height=4.5cm, keepaspectratio]{25-24.697}
	\end{minipage}
	\hspace{-20mm}%
	\begin{minipage}{0.3\textwidth}
		\centering
		\includegraphics[width=\textwidth, height=4.5cm, keepaspectratio]{25-24.696}
	\end{minipage}
	\hspace{-20mm}%
	\begin{minipage}{0.3\textwidth}
		\centering
		\includegraphics[width=\textwidth, height=4.5cm, keepaspectratio]{25-24.681}
	\end{minipage}
	\caption{Amulets depicting a winged Pataikos, a winged ram-headed god and a ram-headed god (MFA 24.697, 24.696, 24.681).}
	\label{fig:25-royal-pataikos}
\end{figure}


An unnamed queen's tomb (Ku. 55) contained an amulet featuring the head of the goddess Hathor on top of a crystal ball. The ball is pierced vertically by a gold tube into a gold base (Figure \ref{fig:25-hathor}). The Napatan kings worshiped Hathor, who was associated with the sacred mountain Jebel Barkal and the \textit{uraeus}. They dedicated a temple to her there (B 200) under her combined form as the goddess Hathor-Tefnut \parencites[.~248]{lohwasserKhonsuSittingJebel2015}[p.~459]{kendallJebelBarkalKarnak2020}.


\begin{figure}[H]
	\centering
	\begin{minipage}{0.45\textwidth}
		\centering
		\includegraphics[width=\textwidth, height=5cm, keepaspectratio]{25-21.321}
		\caption{Amulet featuring a Hathor head on top of a crystal ball (MFA 21.321).}
		\label{fig:25-hathor}
	\end{minipage}%
	\hspace{10pt}
	\begin{minipage}{0.45\textwidth}
		\centering
		\includegraphics[width=\textwidth, height=5cm, keepaspectratio]{25-24.682}
		\caption{Amulet featuring a Hathor head on top of a crystal ball (MFA 21.321).}
		\label{fig:25-2heads}
	\end{minipage}
\end{figure}


Another unusual and unique amulet found with an unidentified queen of Piankhy (Ku. 51) features a deity with two hawk heads, each one with the sun disk and \textit{uraeus} (Figure \ref{fig:25-2heads}). It is possible that this figure represents the god Khonsu, as he could be portrayed with a human body and two hawk or falcon heads \parencites{kniggesalisKhonsu2009, adelzakinasrApotropaicRolesKhonsu2022}. However, this representation does not appear to be consistent in Egypt and few sources feature this type of iconography. This amulet appears to be one of a kind, without parallels in Egypt or Nubia. In Napata, Khonsu was worshiped as the son of the syncretic, lion-headed Mut-Hathot-Tefnut \parencites[p.~353]{gaboldeAmunCultIts2020}. An amulet from Queen Tabiry's tomb (Ku. 53) features a representation of Ra and Amun side-by-side (Figure \ref{fig:25-ra-amun}). Given the similarities between the Ra figure with the double hawk-headed deity amulet, it is also possible that this particular amulet is actually a representation of Ra with two heads. Furthermore, the depiction of a double-headed hawk is known from the Kerma culture, although they lack the distinctive eye markings of Egyptian deities \parencites[p.~173]{minorUseEgyptianEgyptianizing2012}.


From Queen Neferukakashta's burial (Ku. 52) comes an amulet depicting a goddess nursing a queen (Figure \ref{fig:25-nursing-queen}). This silver amulet is particularly interesting as it is, in concept, similar to the Isis nursing Horus amulets. However, instead of an infant Horus seated on his mother's lap, the goddess breastfeeds a queen while both stand up. The queen wears the Kushite cap crown with one \textit{uraeus}, where a vulture rests, further highlighting the combination of Kushite and Egyptian elements \parencites{törökRoyalCrownsKush1987, labiaGoddessQueenGoddess2023}.


\begin{figure}[H]
	\centering
	\begin{minipage}{0.45\textwidth}
		\centering
		\includegraphics[width=\textwidth, height=5cm, keepaspectratio]{25-24.625}
		\caption{Amulet of Ra and Amun from Queen Tabiry's tomb (Ku. 53) (MFA 24.625).}
		\label{fig:25-ra-amun}
	\end{minipage}%
	\hspace{10pt}
	\begin{minipage}{0.45\textwidth}
		\centering
		\includegraphics[width=\textwidth, height=5cm, keepaspectratio]{25-24.928}
		\caption{Amulet of Isis nursing a queen (MFA 24.928).}
		\label{fig:25-nursing-queen}
	\end{minipage}
\end{figure}


Another amulet from the same tomb has the shape of a plaque with a group of deities in profile, three hawk-headed and two female, framed by a winged sun disc and two \textit{uraei} on top and a row of seven \textit{udjat} at the bottom. In the reverse, there is an incised giraffe feeding on a palm tree and other unidentified incised figures, with some potentially representing feathers (Figure \ref{fig:25-plaques}, c, d). Depictions of giraffes (an animal native to this region) in Nubian artifacts are known since the A-Group, however they were particularly popular in the Kingdom of Kerma. During the New Kingdom occupation, this motif disappears and only returns in the Meroitic period \parencites[p.~189-190]{kilroeGiantsSandsGiraffe2023}. As the only depiction of a giraffe in all Napatan royal burials, this plaque is a unique indicator of its Nubian context.


Ku. 51 also yielded two plaque amulets showcasing multiple deities together (Figure \ref{fig:25-plaques}, a, b). The first features anthropomorphic Amun between Isis and Horus and, on the reverse, Bes surrounded by hieroglyphic signs, an \textit{uraeus}, a crocodile, a lizard and other animals (Figure \ref{fig:25-plaques}, a, b). It is noteworthy that both of these plaque's reverses display motifs of animals commonly found in Nubia, while the front depicts deities from the state religion. I would argue that these particular amulets are good examples for the blend of Egyptian influence and local tradition. Other plaque amulets include one from Ku. 51 of ram-headed Amun along with two falcon-headed deities and, on the back, a \textit{udjat} and winged sun disc; and one from Queen Tabiry (Ku. 53) featuring Amun, Horus and, possibly, Isis, with a winged lion-headed goddess on the reverse. As most of these uncommon amulets come from the same group of Piankhy's queens' tombs (Ku. 51, 52, 53 and 55), it should be considered whether they were restricted to these particular burials during this generation or are a reflection of personal preferences during Piankhy's reign.


In contrast, the few deity amulets at Nuri represent Maat, Hathor and Anubis (Figure \ref{fig:25-royal-deities-nuri}). There are two amulets depicting Maat and one of Hathor in another queen of Tanwetamani' tomb, queen Maloqe (Nu. 59). The tomb of an unidentified minor wife of the same king (Nu. 80) contained one amulet of Anubis. There are only three amulets of Maat in royal burials, all belonging to 25th Dynasty queens, suggesting their use was restricted to this phase and to a particular owner. Although, depictions of Anubis were found in the lunettes of stelae from funerary chapels \parencites[p.~636]{lohwasserNapatanNecropoleisBurial2019}, and, later the god would have a prominent role in offering rituals painted on Meroitic pyramid chapel walls and offering tables \parencites[p.~791-799]{helmbold-doyéTombArchitectureBurial2019}[p.~577]{yellinProlegomenaStudyMeroitic2020}, the examples of his depiction in amulets is very limited.


\begin{figure}[H]
	\centering
	\begin{subfigure}{0.3\textwidth}
		\includegraphics[width=\linewidth, height=4cm, keepaspectratio, clip, trim={2cm 0cm 2cm 0cm}]{25-21.304}
		\caption{Front}
	\end{subfigure}%
	\hspace{-2mm}%
	\begin{subfigure}{0.3\textwidth}
		\includegraphics[width=\linewidth, height=4cm, keepaspectratio]{25-21.304-2}
		\caption{Back}
	\end{subfigure}%
	\hspace{2pt}%
	\begin{subfigure}{0.2\textwidth}
		\includegraphics[width=\linewidth, height=4cm, keepaspectratio]{25-24.656}
		\caption{Front}
	\end{subfigure}%
	\hspace{-3mm}%
	\begin{subfigure}{0.2\textwidth}
		\includegraphics[width=\linewidth, height=4cm, keepaspectratio]{25-24.656-2}
		\caption{Back}
	\end{subfigure}
	\caption{Amulets featuring multiple deities and other motifs (MFA 24.656, MFA 21.304).}
	\label{fig:25-plaques}
\end{figure}


\begin{figure}[H]
 	\centering
 	\includegraphics[width=1\textwidth]{25-royal-deities-nuri}
 	\caption{Distribution of amulets with deity motifs in 25th Dynasty royal tombs by occupant's title at Nuri.}
 	\label{fig:25-royal-deities-nuri}
\end{figure}


In both sites, symbol amulets have a less diverse set of motifs (Figure \ref{fig:25-royal-symbols}). Among them are the two most common across all categories: the \textit{udjat} and the \textit{menat} (Figure \ref{fig:25-symbols-pics}). The \textit{udjat} comprises a total of 168 amulets in El-Kurru and 15 at Nuri. While both kings' and queens' tombs at El-Kurru contained these amulets (123 and 45, respectively), at Nuri, 12 amulets were found in queens' tombs, two in kings' and only one in the burial of a minor wife. The \textit{menat} totals 39 instances at El-Kurru (20 with kings and 19 with queens), while there are 16 at Nuri, exclusively in kings' tombs. In particular, some of the \textit{menat} amulets were found in situ placed on the stairs of the tombs (in the tombs of Ku. 18, Ku. 16 and Ku. 5), while others were found strung together with beads (Ku. 5). Therefore, this amulet appears to have taken an additional use and meaning during the Napatan period and could have been placed on the stairs after the closing of the tomb \parencites[p.~246]{howleyRoyalPyramidTombs2015}. Other amulets common to both sites are the Isis knot (\textit{tjet}) and the heart (\textit{ib}) (Figure \ref{fig:25-symbols-pics}).

El-Kurru also features unique symbol amulets. One of those is the winged scarab, which appears exclusively in the tombs of queens that already included some of the winged deity motifs (Ku. 51, 52, 53, 54 and 55) (Figure \ref{fig:25-scarabs}).


\begin{figure}[H]
	\centering
	\begin{minipage}{1\textwidth}
		\includegraphics[width=\textwidth, height=10cm, keepaspectratio]{25-royal-symbols-el-kurru}
	\end{minipage}
	
	\begin{minipage}{1\textwidth}
		\includegraphics[width=\textwidth, height=8cm, keepaspectratio]{25-royal-symbols-nuri}
	\end{minipage}
	\caption{Distribution of amulets with symbol motifs in 25th Dynasty royal tombs by occupant's title.}
	\label{fig:25-royal-symbols}
\end{figure}
 
 
\begin{figure}[H]
 	\centering
 	\begin{minipage}{0.3\textwidth}
 		\includegraphics[width=\textwidth, height=4cm, keepaspectratio]{25-24.650}
 	\end{minipage}
 	\hspace{-2mm}%
 	\begin{minipage}{0.25\textwidth}
 		\includegraphics[width=\textwidth, height=4cm, keepaspectratio]{25-24.586}
 	\end{minipage}
 	\hspace{-15mm}%
 	\begin{minipage}{0.25\textwidth}
 		\includegraphics[width=\textwidth, height=4cm, keepaspectratio]{25-24.1086}
 	\end{minipage}
 	\hspace{-15mm}%
 	\begin{minipage}{0.3\textwidth}
 		\includegraphics[width=\textwidth, height=5cm, keepaspectratio]{25-20.648}
 	\end{minipage}%
 	\caption{Amulets depicting the \textit{udjat}, \textit{menat}, \textit{tjet} and ib from Ku. 53 and Nu. 59 (MFA 24.650, 24.586, 24.1086 and 20.648).}
 	\label{fig:25-symbols-pics}
\end{figure}
 
 
 
Two unique ram-headed amulets were found in the tomb of an unnamed queen of Piankhy (Ku. 55) (Figure \ref{fig:25-ram-sphinx}). They depict the criosphinx, a creature with the head of a ram and the body of a winged lion. They were associated with the god Amun. The creatures are seated on top of columns made of different stones. The motif does not seem to be very common, as the two amulets occur only at El-Kurru during this phase and are different in design, indicating they were not mass manufactured. They bare similarities with the amulet depicting a cat on top of a papyrus column \parencite[p.~33]{andrewsAmuletsAncientEgypt1994}.


Other amulets unique from El-Kurru, each represented by a single example, depict a seated sphinx and the hieroglyphic signs \textit{nt}--- associated with the goddess Neith---, and \textit{nb}---meaning gold (Figure \ref{fig:25-symbols-pics2}). Nuri's exclusive amulets include those of a \textit{ba} bird, \textit{djed} pillars, \textit{uraei}, \textit{nefer} (hieroglyphic symbol for “perfect”) and double feather (Figure \ref{fig:25-symbols-pics3}). \textit{djed} pillar amulets, in particular, have been found in Kerman burials \parencites[p.~145]{minorUseEgyptianEgyptianizing2012}. The \textit{uraeus}, a royal symbol, comprised four examples and were all found in the tomb of queen Malaqaye (Ku. 59).
 
 
\begin{figure}[H]
	\centering
	\begin{minipage}{0.45\textwidth}
		\begin{subfigure}{0.48\textwidth}
			\centering
			\includegraphics[width=\textwidth, height=4cm, keepaspectratio]{25-24.662}
		\end{subfigure}
		\hspace{-2mm}%
		\begin{subfigure}{0.48\textwidth}
			\includegraphics[width=\textwidth, height=4cm, keepaspectratio]{25-24.713}
		\end{subfigure}
		\caption{Winged scarab amulets from Ku. 53 (MFA 24.662 and 24.713).}
		\label{fig:25-scarabs}
	\end{minipage}%
	\hspace{10pt}
	\begin{minipage}{0.45\textwidth}
		\begin{subfigure}{0.48\textwidth}
			\centering
			\includegraphics[width=\textwidth, height=4cm, keepaspectratio]{25-24.972}
		\end{subfigure}
		\hspace{-8mm}%
		\begin{subfigure}{0.48\textwidth}
			\includegraphics[width=\textwidth, height=4cm, keepaspectratio]{25-24.974}
		\end{subfigure}
		\caption{Amulets depicting a ram head sphinx on top of a column (MFA 24.972, 24.974).}
		\label{fig:25-ram-sphinx}
	\end{minipage}
\end{figure}


\begin{figure}[H]
	\centering
	\begin{minipage}{0.25\textwidth}
		\includegraphics[width=\textwidth, height=4cm, keepaspectratio]{25-24.630}
	\end{minipage}
	\hspace{-18mm}%
	\begin{minipage}{0.25\textwidth}
		\includegraphics[width=\textwidth, height=4cm, keepaspectratio]{25-24.1014}
	\end{minipage}
	\hspace{-18mm}%
	\begin{minipage}{0.35\textwidth}
		\includegraphics[width=\textwidth, height=4cm, keepaspectratio]{25-24.1015}
	\end{minipage}%
	\caption{Amulets of a sphinx, nt and neb (MFA 24.630, 24.1014 and 24.1015).}
	\label{fig:25-symbols-pics2}
\end{figure}


\begin{figure}[H]
	\centering
	\begin{minipage}{0.25\textwidth}
		\includegraphics[width=\textwidth, height=4cm, keepaspectratio]{25-20.670}
	\end{minipage}
	\hspace{-3mm}%
	\begin{minipage}{0.20\textwidth}
		\includegraphics[width=\textwidth, height=4cm, keepaspectratio]{25-20.673}
	\end{minipage}
	\hspace{-10mm}%
	\begin{minipage}{0.20\textwidth}
		\includegraphics[width=\textwidth, height=4cm, keepaspectratio]{25-20.671}
	\end{minipage}
	\hspace{-10mm}%
	\begin{minipage}{0.20\textwidth}
		\includegraphics[width=\textwidth, height=4cm, keepaspectratio]{25-20.654}
	\end{minipage}
	\hspace{-10mm}%
	\begin{minipage}{0.20\textwidth}
		\includegraphics[width=\textwidth, height=4cm, keepaspectratio]{25-20.247}
	\end{minipage}%
	\caption{Amulets depicting a \textit{ba} bird, \textit{djed} pillar, \textit{uraeus}, \textit{nefer} and double feather (MFA 20.670, 20.673, 20.671, 20.654 and 20.247).}
	\label{fig:25-symbols-pics3}
\end{figure}


\begin{figure}[H]
	\centering
	\includegraphics[width=1\textwidth]{25-royal-animals}
	\caption{Distribution of amulets with animal motifs in 25th Dynasty royal tombs by occupant's title.}
	\label{fig:25-royal-animals}
\end{figure}


Regarding animal amulets, El-Kurru contained a less diverse assemblages than previous motifs (Figure \ref{fig:25-royal-animals}). Ram, hawk and vulture amulets are common to both sites, with the frog only appearing at Nuri. A few animal amulets were paired with symbols. For instance, the ibis amulet also featured a Maat feather, while a falcon with drooping wings had the sun disk on its head. Additionally, all cow amulets from Nuri featured the Hathor headdress (Figure \ref{fig:25-animals-pics}).



\begin{figure}[H]
	\centering
	\begin{minipage}{0.3\textwidth}
		\includegraphics[width=\textwidth, height=3.5cm, keepaspectratio]{25-24.1085}
	\end{minipage}
	\hspace{-9.5mm}%
	\begin{minipage}{0.31\textwidth}
		\includegraphics[width=\textwidth, height=4cm, keepaspectratio]{25-24.973}
	\end{minipage}
	\hspace{-8mm}%
	\begin{minipage}{0.3\textwidth}
		\includegraphics[width=\textwidth, height=4cm, keepaspectratio]{25-20.669}
	\end{minipage}

	\begin{minipage}{0.3\textwidth}
		\includegraphics[width=\textwidth, height=4cm, keepaspectratio]{25-24.706}
	\end{minipage}
	\hspace{-22mm}%
	\begin{minipage}{0.3\textwidth}
		\includegraphics[width=\textwidth, height=4cm, keepaspectratio]{25-21.314}
	\end{minipage}
	\hspace{-13mm}%
	\begin{minipage}{0.3\textwidth}
		\includegraphics[width=\textwidth, height=4cm, keepaspectratio]{25-24.976}
	\end{minipage}%
	\caption{Amulets of a cow, ibis, hawk, falcon, fly and ram (MFA 24.1085, 24.973, 20.669, 24.706, 21.314, 24.976).}
	\label{fig:25-animals-pics}
\end{figure}


Another unique amulet was found at El-Kurru: it features a silver ram's head on top of a crystal pillar (Figure \ref{fig:25-animal-pics}, lower lower right). This amulet comes from the tomb of an unidentified queen of Piankhy (Ku. 55), which also contained the other ram-headed sphinx amulets. Taharqa's tomb also features another unique ram amulet with the sun disk and double \textit{uraeus}. Dunham described the gold amulet as a ram's head “soldered to a granulated cylinder with two loops” \parencites[p.~12]{dunhamNuri1955}. Unfortunately, I have not been able to uncover the current location of the amulet in order to examine it in more detail and include a photograph. Therefore, we only have the drawing from Dunham's report (Figure \ref{fig:25-ram-pillar}). Unlike the El-Kurru amulet, where the pillar is positioned under the ram's head, the one from Nuri has the cylinder attached to the top of the ram's head.


\begin{wrapfigure}{l}{0.2\textwidth}
	\centering
	\includegraphics[width=1\linewidth, height=20cm, keepaspectratio]{25-ram-pillar}
	\caption{Drawing of a ram's head amulet attached to a “cylinder” \parencites[p.~15, fig.~5]{dunhamNuri1955}}
	\label{fig:25-ram-pillar}
\end{wrapfigure}


El-Kurru contains the only fly amulet of the royal sites. In Nubia, fly amulets have been found in burials and urban contexts dated to the Kerma period, as well as New Kingdom fortresses at Semna and houses at Aniba \parencites{ditriaAmuletsKermaCulture2022, sidpuraFliesLionsOyster2022}. In Egypt, fly amulets have been found in the tombs of New Kingdom queens. In fact, the fly amulets in Queen Anhotep's tomb has parallels with the Kerman amulets, suggesting the adoption of this motif by the Egyptian elite \parencites{lilyquistTombThreeForeign2003, lacovaraTreasureEgyptianQueen2022}. Still, it appears that this amulet motif was not popular in Napatan royal burials. 

Overall, the animal amulets contained both animals associated with certain Egyptian deities---such as cows with Hathor, hawks with Horus, ibis with Thoth and cats with Bastet and Sekhmet---, others could be more representative of the natural environment, namely monkeys, apes, fishes and frogs, although most of them comprised of a single example.


\begin{figure}[H]
	\centering
	\includegraphics[width=1\textwidth]{25-royal-nature-objs-misc}
	\caption{Distribution of amulets with nature, object and miscellaneous motifs in 25th Dynasty royal tombs by the occupant's title.}
	\label{fig:25-royal-nature-objs-misc}
\end{figure}


Since amulets depicting nature elements, objects and miscellaneous motifs are very rare, they were all included in the same graph (Figure \ref{fig:25-royal-nature-objs-misc}). All were found in the tombs of queens. The most common motif in these categories is the papyrus column, the only amulet to occur both in El-Kurru and Nuri (four and six examples, respectively). Amulets depicting objects include a headrest from Nuri, and a vase from El-Kurru (Figure \ref{fig:25-objs-pics}, top row).

The miscellaneous/other category include only two amulets. One is an incised cylinder depicting a vulture wearing the sun disk and featuring some hieroglyphic symbols (Figure \ref{fig:25-objs-pics}, third). This particular amulet is more similar to a bead, however it was included in the database due to the decoration and inscriptions. Also, it was not considered an amulet of a vulture due to the shape not being that of the animal. The second amulet represents an unknown motif. In Dunham's report, there is no indication of the motif and, unfortunately, the quality of the photograph did not allow me to make a positive identification \parencites[p.~101, pl.~XX.A, g, 3]{dunhamElKurru1950}. Additionally, the location of this amulet is unknown, complicating its identification \parencites{hinkelCatalogueObjectsSudan2002, kormyshevaGodsDivineSymbols2006}.



 \begin{figure}[H]
	\centering
	\begin{minipage}{0.25\textwidth}
		\includegraphics[width=\textwidth, height=4cm, keepaspectratio]{25-24.651}
	\end{minipage}
	\hspace{-13mm}%
	\begin{minipage}{0.25\textwidth}
		\includegraphics[width=\textwidth, height=4cm, keepaspectratio]{25-20.649}
	\end{minipage}
	\hspace{-13mm}%
	\begin{minipage}{0.25\textwidth}
		\includegraphics[width=\textwidth, height=4cm, keepaspectratio]{25-24.617}
	\end{minipage}
	\hspace{-18mm}%
	\begin{minipage}{0.35\textwidth}
		\includegraphics[width=\textwidth, height=4cm, keepaspectratio]{25-20.650}
	\end{minipage}%
	\caption{Amulets of a vase, headrest, papyrus column and incised cylinder (MFA 24.651, 20.650, 20.649, 24.617).}
	\label{fig:25-objs-pics}
\end{figure}
 
 
\subsubsection{Amulet materials}


\begin{figure}[H]
	\centering
	\begin{minipage}{1\textwidth}
		\centering
		\includegraphics[width=\textwidth, height=10cm, keepaspectratio]{25-royal-materials-el-kurru}
	\end{minipage}
	\begin{minipage}{1\textwidth}
		\centering
		\includegraphics[width=0.92\textwidth, height=10cm, keepaspectratio]{25-royal-materials-nuri}
	\end{minipage}
	\caption{Distribution of amulet materials in 25th Dynasty royal tombs by occupant's title.}
	\label{fig:25-royal-materials}
\end{figure}


There are differences in materials between El-Kurru and Nuri---although faience is the most common material in both sites (Figure \ref{fig:25-royal-materials}). At El-Kurru, the majority of faience amulets come from queens, while at Nuri, they come predominantly from kings' tombs. It is not surprising that faience was such a popular material for amulets---excavations at the Amun Temple in nearby Sanam have uncovered evidence of faience production, indicating that objects made from this material were likely manufactured locally for the royal burials \parencites{reisnerNoteHarvardBostonExcavations1920, howleyNewDiscoveriesSanam2020}. At both sites, amulets of queens were made from the most varied materials, followed by kings. Amulets of minor wives at El-Kurru were only made of faience, and, at Nuri, they were composed of gold, beryl and serpentinite.


Amulets at El-Kurru have more materials made from foreign materials, mostly sourced at the central part of the Eastern Desert. Among them, jasper, serpentinite and travertine \parencites[p.~30, 503-506, 790]{harrellArchaeologyGeologyAncient2024}. As during the 25th Dynasty, the kings of Kush controlled part of the Egyptian territory, access to these materials would be facilitate. Regarding jasper, present in both sites, depending on the coloring, the stone could have come from different places. In the reports, Dunham either specifies a red coloration or not at all. Analysis of the jasper amulets' registries at the MFA's online catalogue revealed that, of the three jasper amulets, one was noted as being red. Sources of red jasper are more abundant in the Egyptian Eastern Desert and there does not appear to be any in the Sudan \parencites[p.~790-791]{harrellArchaeologyGeologyAncient2024}.


An imported material from further away is lapis lazuli, which is the second most common material for amulets at El-Kurru, without instances at Nuri. Lapis lazuli was sourced from modern-day Afghanistan, as this is the only known source of this material in antiquity \parencites[p.~869]{harrellArchaeologyGeologyAncient2024}. Some materials are more difficult to assess if they were imported or local, such as silver, which was used for seventeen amulets from El-Kurru. Silver, for example, could have been mined from Egyptian Eastern Desert deposits after the 18th Dynasty \parencites[p.~878]{harrellArchaeologyGeologyAncient2024}---although, analysis of gilded silver objects from Nuri at the MFA could not determine the source of the material \parencites{gänsickeGildedSilverAncient2000}. Further analysis of silver objects from Third Intermediate and Late Period from Egypt also did not result in a definite identification of the source. The author also considers the possibility of reuse and recycling of old silver stored at temples \parencites[p.~59-60]{jurmanSilverTreasuryHerishef2015}. Therefore, at the moment, the source of silver during the Napatan period cannot be determined without further analysis. Even so, the present evidence suggests it either came from Egypt or was imported from other locations.


Beryl is an interesting case, as was found at both El-Kurru and Nuri. Beryl is a gemstone with two color varieties, each with their own denomination: aquamarine for blue and emerald for green \parencites[p.~758]{harrellArchaeologyGeologyAncient2024}. According to Dunham's report on Nuri, "an attractive greenish semiprecious stone" was mostly recorded as beryl \parencites[p.~5]{dunhamNuri1955}. However, Harrell warns that non-geologists have misidentified amazonite used in the 18th and 19th Dynasties in Egypt as green beryl/emerald and that beryl was actually used only during the Ptolemaic and Roman periods \parencites[p.~758, 807]{jurmanSilverTreasuryHerishef2015}. Therefore, it is possible that the material Reisner/Dunham identified as beryl is actually amazonite. According to Harrell, there are only three known ancient mines for amazonite: two in Egypt, both in the Eastern Desert region of Thebes and one near the Egyptian-Sudanese border \parencites[p.~795-802]{harrellArchaeologyGeologyAncient2024}. Analysis conducted on amazonite beads from a Neolithic site in northern Sudan confirmed that the material came from a deposit in southern east Ethiopia \parencites[p.~1202]{zerboniLongdistanceExchangeAmazonite2018}. With this evidence, we can speculate where the amazonite from the royal sites originated. On one hand, it is possible that the material was sourced in mines between Thebes and Lower Nubia. On the other hand, trade routes between Kush and Ethiopia could have been maintained. A further consideration is the Land of Punt---speculated to be in present-day Eritrea \parencites{grathwolAdulisTransshipmentBaboons2023}---, where Egyptians obtained various imported goods.


Materials restricted to a few amulets include locally available ones, such as marble, electrum, malachite and ivory \parencites[p.~927-929]{harrellArchaeologyGeologyAncient2024}. The unique ivory amulet is intriguing, as this was an exported commodity from Nubia and even that 163 kg of unworked ivory fragments were discovered at the Treasury storerooms in Sanam \parencites{griffithOxfordExcavationsNubia1922}[p.~128-130]{vincentelliLongDistanceTrade2018}{shochatThousandYearsNubian2025}. Such evidence suggests that ivory may not have been a preferred material specifically for amulets.


The unique enamel amulet at El-Kurru depicts a falcon with the double crown and is described by Dunham as “cloisonné enamel” (Figure \ref{fig:25-falcon-enamel}) \parencites[p.~93]{dunhamElKurru1950}. However, this is more related to the technique than the material. Cloisonné enamel describes the process of decorating metalwork with coloured glass. In Egypt, only two objects from the Third Intermediate Period, specifically 21st Dynasty, are considered to be made of true enamel. The technique was, however, employed during the Meroitic period \parencites[p.~166]{ogdenMetals2000}. Unfortunately, Dunham's report did not specify which metal the Napatan amulet was made of and I have been unable to locate its current whereabouts. It is possible that one of the amulets depicting a ram-headed sphinx on top of a column also used the same technique, as its description the MFA's online catalogue mentions glass inlays. Should the falcon amulet be found, analysis of both amulets could confirm whether glass enamel technology was employed during the Napatan period, even if only in exceptional cases.


\begin{wrapfigure}{l}{0.2\textwidth}
	\centering
	\includegraphics[width=1\linewidth, height=20cm, keepaspectratio]{25-falcon-enamel}
	\caption{Amulet of a falcon made of metal and described as “ cloisonné enamel” \parencites[pl.~LX.D, I]{dunhamNuri1955}}
	\label{fig:25-falcon-enamel}
\end{wrapfigure}


The amulet assemblage at Nuri features exclusive materials, such as rock crystal, quartz, haematite, carnelian, chalcedony and agate. Provenance analysis suggests a reliance on sources from the Eastern Desert for many of these stones \parencites[p.~759]{harrellArchaeologyGeologyAncient2024}. The material with the most amulets was rock crystal, sourced from amethyst mines of the Eastern Desert in Lower Nubia. All recovered examples were \textit{udjat} amulets, found exclusively in the tomb of Tanwetamani's queen, Malaqaye (Nu. 59). Milky quartz, also likely sourced from Aswan, was described in the reports as just quartz or “white quartz”. Analysis of images from the MFA confirms that the amulets were, in fact, made of milky quartz and were fashioned into symbol motifs (\textit{uraeus}, heart and \textit{djed}). Haematite was a gemstone available from a mine in Northern Sudan, while carnelian and chalcedony could be sourced from Stela Ridge, south of Aswan, and agate could be found in Nile River pebbles or mines in the Sinai and Eastern Desert  \parencites[p.~759, 866, 948]{harrellArchaeologyGeologyAncient2024}. In particular, carnelian has been used for amulets in Nubia since the Kerma culture \parencites[p.~146]{minorUseEgyptianEgyptianizing2012}. However, the use of these materials for amulets was exceptionally rare during the 25th Dynasty.


In summary, amulets could have been made from a variety of materials, including imported and local ones. However, faience was the most common material, which was likely manufactured in nearby Sanam. Evidence of trade between the Kushites and other lands come from the Treasury and workshops in Sanam containing several materials and objects, which was active from the reign of Piankhy until Aspelta \parencites{vincentelliLongDistanceTrade2018}. Excavations revealed the presence of carnelian, lapis lazuli, agate, crystal and marble; foreign and local materials found in the amulets. This discovery uncovered trade routes connecting Kush to different lands, including the Near East, Egypt and southern lands. Furthermore, a workshop contained working tools for stone and faience, including a mould for a \textit{udjat} and polishing material. Therefore, it is possible that materials arriving at Sanam were used for the manufacture of objects, including amulets.


\section{Early Napatan}


\subsection{Tomb structure}


\begin{figure}[H]
	\centering
	\includegraphics[width=1\textwidth]{EN-royal-tombs}
	\caption{Distribution of superstructure and substructure types by occupant's title in Early Napatan royal burials.}
	\label{fig:EN-royal-tombs}
\end{figure}


Following Reisner's dating, the tomb of an unnamed queen of Atlanersa (Ku. 61) stands out as the only Early Napatan burial at El-Kurru, while all other tombs were built at Nuri (Figure \ref{fig:EN-royal-tombs}). However, it is possible that this tomb is actually from the late 25th Dynasty and was misdated by Reisner. These tombs maintain the 25th Dynasty's established substructure of underground chambers. Regarding superstructures, the minor wives' tombs are similar to those from the pre- and 25th Dynasty phases, as they are all unidentified. Queens' tombs mostly exhibit pyramid superstructures, with some unknown ones, whereas all kings tombs were pyramids.


Chapels were present in the majority of royal tombs of this phase---with the exception of one queen (Nu. 58) and all tombs of minor wives (Table \ref{tab:EN-tomb-elements}). Stairs are the only form of entrance, as established during the 25th Dynasty. The chamber configuration in Early Napatan royal burials reveal patterns, as well as some exceptions. Beginning with Senkamanisken's tomb (Nu. 3), all kings' burials feature three chambers, except Atlanersa's (Nu. 20), the earliest tomb of this phase, which had two chambers. Minor wives' tombs consistently contained a single chamber. Queens' tombs, however, show slight variations. Before Senkamanisken's burial, out of three queens' tombs, two had one chamber (Nu. 53 and 60) and one had two (Ku. 61). After Senkamanisken's burial, most queens' tombs include two chambers, with sole exception of queen Artaha's tomb (Nu. 58), which included one chamber.


Decoration in chapels and chambers seems to be linked to title. While the majority of king's and queen's tombs were decorated, none of the minor wives' were. As always, preservation issues could account for their absence. An interesting hypothesis was posited regarding the choice of decoration for queen Yeturow's chamber (Nu. 53): as neither Reisner's early excavations nor recent work by the Nuri Archaeological Expedition (NAE) uncovered traces of a chapel, researchers proposed that its absence could explain the choice of wall decoration in this particular tomb, with mummiform figures requesting food offerings and texts related to sustenance \parencites[p.~201]{strongReexaminingTombQueen2021}.


The sole tomb at El-Kurru (Ku. 61) is is the only one to include both bed cutouts and a coffin bench. As for the Nuri tombs, holes for a bed were completely absent, while the inclusion of coffin benches is inconsistent and it does not appear to follow a pattern regarding the title of the owner. Nevertheless, two tombs contained stone sarcophagi, those of kings Anlamani (Nu. 6) and Aspelta (Nu. 8). The sarcophagi's decoration and texts originated from different Egyptian sources, such as the Coffin and Pyramid texts, and the Book of the Dead. The design of Anlamani's sarcophagus was unique both in Kush and Egypt, and was eventually copied by Aspelta \parencites{dollTextsDecorationNapatan1978, dollIdentificationSignificanceTexts1982, dollRoyalMortuaryCult2014}. Aspelta's tomb featured a coffin bench, whereas Anlamani's did not. This may suggest that the use of a sarcophagus was not connected with the inclusion of a bench. The presence of surviving coffin fittings and eye inlays in some tombs suggest the use of wooden containers that, unfortunately, did not survive.


{\small
	\begin{xltabular}{\textwidth}{X X X X X X X}
		\caption{Distribution of architectural elements by tomb owner's title in Early Napatan royal burials.}
		\label{tab:EN-tomb-elements} \\
		\toprule
		\textbf{Tomb} & \textbf{Owner} & \textbf{Chapel} & \textbf{Type of entrance} & \textbf{N. of chambers} & \textbf{Decoration} & \textbf{Bed holes and/or coffin bench} \\
		\midrule
		\endfirsthead
		\multicolumn{7}{c}{\textit{Continued from previous page}} \\
		\toprule
		\textbf{Tomb} & \textbf{Owner} & \textbf{Chapel} & \textbf{Type of entrance} & \textbf{N. of chambers} & \textbf{Decoration} & \textbf{Bed holes and/or coffin bench} \\
		\midrule
		\endhead
		\bottomrule
		\endlastfoot
		Nu. 20 & king & yes & stairs & 2 & -- & coffin bench \\
		\cmidrule(r){1-7}
		Ku. 61 & queen & -- & stairs & 2 & yes & both \\
		\cmidrule(r){1-7}
		Nu. 53 & queen & -- & stairs & 1 & yes & coffin bench \\
		\cmidrule(r){1-7}
		Nu. 60 & queen & -- & stairs & 1 & -- & coffin bench \\
		\cmidrule(r){1-7}
		Nu. 75 & minor wife & -- & stairs & 1 & -- & -- \\
		\cmidrule(r){1-7}
		Nu. 3 & king & yes & stairs & 3 & yes & coffin bench \\
		\cmidrule(r){1-7}
		Nu. 41 & queen & yes & stairs & 2 & -- & coffin bench \\
		\cmidrule(r){1-7}
		Nu. 71 & minor wife & -- & stairs & 1 & -- & -- \\
		\cmidrule(r){1-7}
		Nu. 73 & minor wife & -- & stairs & 1 & -- & coffin bench \\
		\cmidrule(r){1-7}
		Nu. 78 & minor wife & -- & stairs & 1 & -- & -- \\
		\cmidrule(r){1-7}
		Nu. 81 & minor wife & -- & stairs & 1 & -- & -- \\
		\cmidrule(r){1-7}
		Nu. 82 & minor wife & -- & stairs & 1 & -- & -- \\
		\cmidrule(r){1-7}
		Nu. 6 & king & yes & stairs & 3 & yes & -- \\
		\cmidrule(r){1-7}
		Nu. 21 & queen & yes & stairs & 2 & yes & coffin bench \\
		\cmidrule(r){1-7}
		Nu. 22 & queen & yes & stairs & 2 & -- & coffin bench \\
		\cmidrule(r){1-7}
		Nu. 23 & queen & yes & stairs & 2 & yes & -- \\
		\cmidrule(r){1-7}
		Nu. 72 & minor wife & -- & stairs & 1 & -- & -- \\
		\cmidrule(r){1-7}
		Nu. 76 & minor wife & -- & stairs & 1 & -- & -- \\
		\cmidrule(r){1-7}
		Nu. 79 & minor wife & -- & stairs & 1 & -- & -- \\
		\cmidrule(r){1-7}
		Nu. 8 & king & yes & stairs & 3 & yes & coffin bench \\
		\cmidrule(r){1-7}
		Nu. 24 & queen & yes & stairs & 2 & yes & -- \\
		\cmidrule(r){1-7}
		Nu. 27 & queen & yes & stairs & 2 & -- & -- \\
		\cmidrule(r){1-7}
		Nu. 40 & queen & yes & stairs & 2 & -- & -- \\
		\cmidrule(r){1-7}
		Nu. 42 & queen & yes & stairs & 2 & -- & -- \\
		\cmidrule(r){1-7}
		Nu. 58 & queen & -- & stairs & 1 & yes & -- \\
	\end{xltabular}
}


\subsection{Object types}


\begin{figure}[H]
	\begin{subfigure}{0.37\textwidth}
		\raggedleft
		\includegraphics[width=\linewidth, height=12cm, keepaspectratio]{EN-royal-objs-el-kurru}
	\end{subfigure}
	\hspace{-2mm}%
	\begin{subfigure}{0.6\textwidth}
		\raggedright
		\includegraphics[width=\linewidth, height=12cm, keepaspectratio]{EN-royal-objs-nuri}
	\end{subfigure}
	\caption{Distribution of object types in Early Napatan royal tombs by occupant's title.}
	\label{fig:EN-royal-objs}
\end{figure}


Object types at El-Kurru come from the sole queen tomb. The majority of grave goods consist of shabtis, with 231 examples (Figure \ref{fig:EN-royal-objs}). Adornments, beads and burial container trappings comprise only eight, seven and three items, respectively; while mummy trappings and non-ceramic vessels had one item each. At Nuri, there is a bigger diversity of object types, as the number of tombs increase. Shabtis are also the most popular: 2391 with kings, 2140 with queens and 1920 with minor wives. Kings' tombs contained the most amount of beads, 1060, followed by minor wives' with 161 and queens' with 121. Ceramic vessels were also present in all burials, with queens' containing 167, kings' 68 and minor wives' 13. The other object types have much fewer items in total. Unlike 25th Dynasty tombs, few objects from the Early Napatan phase were exclusive to one type of burial: offering vessels and stelae were only attested with kings, while the only bead net was found with a queen. However, it is possible that individual beads in kings and minor wives' tombs were originally a bead net. Minor wives' continue to be the burials with the least diversity of object types. Furthermore, jar stoppers and weapons---present in 25th Dynasty assemblages---are absent in Early Napatan tombs.
Heart scarabs were found in the three types of burials: one with Aspelta (Nu. 8), one with a minor wife of Senkamanisken (Nu. 81), and three with queens. While Asata, queen of Aspelta (Nu. 42) only contained one heart scarab, queen Yeturow of Atlanersa (Nu. 53) was buried with two. Whether Yeturow was a special case or looting is to be blamed for their absent in other tombs is unknown.
The two stelae came from the burials of Anlamani and Aspelta. Analysis revealed the use of texts from various Egyptian sources, which were intentionally selected (with certain excerpts rejected) and compiled to form an unique Napatan funerary textual composition \parencites{dollRoyalMortuaryCult2014}.


A new type of object, the cylinder sheath, emerged in the Early Napatan phase. These items, crafted in gold and gilded silver, were decorated with motifs of deities---specially winged goddesses---, ram's heads and \textit{uraei}, among others. There were some variations in decoration, albeit they appear to have been made in pairs. The manufacture was made in several stages with elaborate techniques and likely produced locally, as their motifs elude to the Jebel Barkal mythology. These objects were open in one end and could have contained organic material. Their precise function is not clear, although different theories about their function have been proposed \parencites{gänsickeFreshLookCylinder2004, pompeiNapatanCylindricalSheaths2019}. In total,  eighteen were found: fifteen in Aspelta's tomb, three in  Senkamanisken's, and only one with a queen of Aspelta, Madiqen (Nu. 27). It is interesting that these objects were only attested in these particular tombs. This could be due to looting, which would have targeted gold and silver objects. As Anlamani came between these two kings and no cylinder sheaths were found in his tomb or his wives', it does not seem likely that they were only used in Senkamanisken and Aspelta's generations, while entirely skipping Anlamani's.


\subsection{Amulets}


\subsubsection{Amulet distribution}


A total of 189 amulets are distributed across 13 tombs, a sharp decline compared to the 25th Dynasty (Table \ref{tab:EN-royal-distro}). In reality, this is the beginning of a significant decline in the number of amulets of later phases.


All kings' burials contained amulets, as was the case in the 25th Dynasty. However, the median was much higher in this phase (23) than in the 25th Dynasty phase (10), meaning more amulets were found per tomb. Slightly more than half of queens were buried with amulets, whereas only 22\% of minor wives contained these objects. The median of amulets with queen also decline, reaching only 1, while minor wives was maintained at 0.


{\small
	\begin{xltabular}{\textwidth}{X X X X X X}
		\caption{Distribution of amulets in Early Napatan royal tombs by owner's title, including total counts of amulets and tombs, tombs containing amulets, amulet frequency (\%) and median.}
		\label{tab:EN-royal-distro} \\
		\toprule
		\textbf{Owner} & \textbf{Total number of tombs} & \textbf{Total number of amulets} & \textbf{N. of tombs with amulets} & \textbf{\% of tombs with amulets} & \textbf{Median n. of amulets per tomb} \\
		\midrule
		\endfirsthead
		\multicolumn{6}{c}{\textit{Continued from previous page}} \\
		\toprule
		\textbf{Owner} & \textbf{Total number of tombs} & \textbf{Total number of amulets} & \textbf{N. of tombs with amulets} & \textbf{\% of tombs with amulets} & \textbf{Median n. of amulets per tomb} \\
		\midrule
		\endhead
		\bottomrule
		\endlastfoot
		King & 4 & 140 & 4 & 100\% & 23 \\
		\cmidrule(r){1-6}
		Queen & 12 & 39 & 7 & 58\% & 1 \\
		\cmidrule(r){1-6}
		Minor wife & 9 & 10 & 2 & 22\% & 0 \\
		\midrule
		\textbf{Total} & \textbf{25} & \textbf{189} & \textbf{13} & -- & -- \\
	\end{xltabular}
}


\subsubsection{Amulet types}


\begin{figure}[H]
	\centering
	\includegraphics[width=1\textwidth]{EN-royal-types}
	\caption{Distribution of amulet types in Early Napatan royal tombs by occupant's title.}
	\label{fig:EN-royal-types}
\end{figure}


Early Napatan amulets were only found at Nuri (Fig. \ref{fig:EN-royal-types}). Symbols predominate, mirroring the distribution observed in 25th Dynasty royal tombs. Kings were buried with the vast majority of these amulets (138), followed distantly queens (30) and minor wives (4). Deities are the second most common category, though they were far less numerous. The distribution of these amulets was more equitable: four were found in queens' burials, two with minor wives, and one with a king. A similar pattern emerges with animal motifs, which were primarily associated with queens (4) and minor wives (2), with a single example found with a king. The sole object amulet was associated with a minor wife, and the only nature-related amulet was found with a queen.


\subsubsection{Amulet motifs}



By far, the most popular amulet during the Early Napatan phase was the \textit{menat}, with a total of 137 examples (Figure \ref{fig:EN-royal-motifs}). These amulets were found predominantly in kings' burials (117) and less so with queens (20), followig the distribution in the previous phase. The \textit{udjat} is the second most common amulet, albeit with only 21 examples. Of these, seventeen were found with kings, three with queens, and one with a minor wife. The remaining motifs were rare, with fewer than ten examples each; among them, eight are represented by a single amulet.


\begin{figure}[H]
	\centering
	\includegraphics[width=1\textwidth]{EN-royal-motifs}
	\caption{Distribution of amulet motifs in Early Napatan royal tombs by the occupant's title.}
	\label{fig:EN-royal-motifs}
\end{figure}

\begin{wrapfigure}{l}{0.3\textwidth}
	\centering
	\begin{minipage}{\linewidth}
		\centering
		\includegraphics[width=0.45\linewidth, height=5cm, keepaspectratio]{EN-20.381}
		\hspace{-2mm}%
		\includegraphics[width=0.55\linewidth, height=5cm, keepaspectratio]{EN-20.391}
		\caption{Early Napatan deity amulets: Duamutef (MFA 20.381) and Imsety (MFA 20.391).}
		\label{fig:EN-pics}
	\end{minipage}
\end{wrapfigure}

Tombs of minor wives yielded more exclusive amulet motifs than the previous phase. Nevertheless, queens' burials continue to be those with the most diverse corpus of amulet motifs. In fact, the tomb of Asata (Nu. 42), a queen of Aspelta, had the most distinct set of amulets, with nine different motifs, including those of deities, animals, symbols and nature. Among them are two amulets depicting the gods Duamutef and Imsety, part of the Four Sons of Horus and related to mummification. These amulets seem very similar in manufacture and may indicate that the amulets were part of a set (Figure \ref{fig:EN-pics}).


In addition to the \textit{udjat} and \textit{menat}, other amulet motifs from the previous phase continued to appear---such as the headrest, papyrus column, frog, cow, vulture, \textit{uraeus}, heart, \textit{djed} column, and the goddess Hathor (Figure \ref{fig:EN-pics2}). New motifs include the gods Duamutef and Ptah, and a snake amulet (Figure \ref{fig:EN-pics2}, third). Thus, the diversity, along with the total number of amulets, is significantly lower than in the 25th Dynasty. This discrepancy may reflect the Kushite's loss of control of the Egyptian territory, which may have limited both the import and local production of amulets.


 \begin{figure}[H]
	\centering
	\begin{minipage}{0.24\textwidth}
		\includegraphics[width=\textwidth, height=3.5cm, keepaspectratio]{EN-20.380}
	\end{minipage}
	\hspace{-13mm}%
	\begin{minipage}{0.24\textwidth}
		\includegraphics[width=\textwidth, height=3.5cm, keepaspectratio]{EN-20.374}
	\end{minipage}
	\hspace{-13mm}%
	\begin{minipage}{0.24\textwidth}
		\includegraphics[width=\textwidth, height=3.5cm, keepaspectratio]{EN-20.371}
	\end{minipage}
	\hspace{-17mm}%
	\begin{minipage}{0.24\textwidth}
		\includegraphics[width=\textwidth, height=3.5cm, keepaspectratio]{EN-20.389}
	\end{minipage}%
	\caption{Early Napatan gold amulets: \textit{uraeus} (MFA 20.380), \textit{djed}  pillar (MFA 20.374), snake (MFA 20.371) and vulture (MFA 20.389).}
	\label{fig:EN-pics2}
\end{figure}


\subsubsection{Materials}


\begin{figure}[H]
	\centering
	\includegraphics[width=1\textwidth, height=10cm, keepaspectratio]{EN-royal-materials}
	\caption{Distribution of amulet materials in Early Napatan royal tombs by the occupant's title.}
	\label{fig:EN-royal-materials}
\end{figure}


Faience was the most common material for amulets in Early Napatan royal tombs, comprising of over 75\% (Figure \ref{fig:EN-royal-materials}). This prevalence is almost entirely attributable to the 137 \textit{menat} amulets. Faience amulets were found with kings (121) and queens (22). In contrast, gold, the second most common material, were found in the burials of queens (12) and minor wives (7). The complete absence of gold amulets in kings' burials is most likely a result of looting. Some gold amulets were motifs stamped into sheets (Figures \ref{fig:EN-pics}, \ref{fig:EN-pics2}). These types of amulets have been found in Egypt, mainly sheet gold breastplates depicting vultures in the tombs of Thutmosis III's foreign minor wives and in the burial of a 21st Dynasty Priestess of Amun in Thebes as well as the previously mentioned Late Period burial G7652 at Giza \parencites[p.~198, 313]{lilyquistTombThreeForeign2003}[p.~199]{astonBurialAssemblagesDynasty2009}{dermanuelianDigitalGizaGiza2025}. Most Early Napatan amulet materials were already present in 25th Dynasty tombs, with one addition: granodiorite, which could be found in Eastern Desert quarries \parencites[p.~506-508]{harrellArchaeologyGeologyAncient2024}. As the Treasury at Sanam remained in use until the end of the Early Napatan phase, this could account for the presence of a single amazonite amulet (if it indeed was imported). However, as there are only three amulets made from materials coming from Egypt's Eastern Desert---jasper and granodiorite---may indicate less access to trade routes in order to acquire resources.


\section{Middle Napatan}


\subsection{Tomb structure}


\begin{figure}[H]
	\centering
	\includegraphics[width=1\textwidth]{MN-royal-tombs}
	\caption{Distribution of superstructure and substructure types by occupant's title in Middle Napatan royal burials.}
	\label{fig:MN-royal-tombs}
\end{figure}


By the Middle Napatan phase, all royalty was buried at Nuri (Figure \ref{fig:MN-royal-tombs}). Tomb structures diverge from the previous phases in a very interesting way: it is the first phase where the superstructures of the minor wives' tombs are all pyramids. Additionally, the majority of queens' tombs were pyramids as well (11 out of 13). This could be due to a better preservation, or that minor wives and certain queens were permitted to have pyramids. Kings' tombs comprise of pyramids and one double pyramid, that of king Irike-Amanote (Nu. 12), the last king of the Middle Napatan phase. It has been hypothesized that this pyramid was enlarged at a later date and was, originally, a single pyramid---similarly to the two previous double pyramids \parencites[p.~108]{welsbyKingdomKushNapatan1996}. Underground chambers remained the standard substructure.

Furthermore, almost all royal burials featured a chapel, unlike  previous phases (Table \ref{tab:MN-royal-tombs-elements}). Whether this reflects an increased funerary significance of chapels or better preservation is uncertain. The sole exception is Nu. 45, belonging to a queen of Malonaqen (likely Tagtal). As Reisner observed, the pyramid was destroyed by later Christian house walls, suggesting that the chapel may have suffered the same fate and, in fact, all tombs from this phase were built with an attached chapel \parencites[p.~149]{dunhamNuri1955}.


Stairs continued to be the standard type of entrance, while the number of chambers appears to follow a pattern depending on the occupant's title. All kings' tombs in this phase contained three chambers, while most queens' had two---except for Nu. 45, with one chamber. Minor wives' tombs varied, containing either one or two chambers. As in earlier phases, preserved decoration on chapel or chamber walls was inconsistent.


None of the tombs included cutouts for a burial bed and half contained a coffin bench. Reisner observed that king Karkamani's tomb (Nu. 7) featured a platform for a sarcophagus but no traces of stone, leading him to conclude that if there was one, it would likely have been made from wood \parencites[p.~161]{dunhamNuri1955}. Regarding other body containers, out of the 30 total tombs, 20 contained either coffin fittings or eye inlays. Since some of those did not include a coffin bench, it is possible that this feature was not necessary for the use of a coffin, and they were most likely used.


{\small
	\begin{xltabular}{\textwidth}{X X X X X X X}
		\caption{Distribution of architectural elements by tomb owner's title in Middle Napatan royal burials.}
		\label{tab:MN-royal-tombs-elements} \\
		\toprule
		\textbf{Tomb} & \textbf{Owner} & \textbf{Chapel} & \textbf{Type of entrance} & \textbf{N. of chambers} & \textbf{Decoration} & \textbf{Bed holes and/or coffin bench} \\
		\midrule
		\endfirsthead
		\multicolumn{7}{c}{\textit{Continued from previous page}} \\
		\toprule
		\textbf{Tomb} & \textbf{Owner} & \textbf{Chapel} & \textbf{Type of entrance} & \textbf{N. of chambers} & \textbf{Decoration} & \textbf{Bed holes and/or coffin bench} \\
		\midrule
		\endhead
		\bottomrule
		\endlastfoot
		Nu. 9 & king & yes & stairs & 3 & yes & -- \\
		\cmidrule(r){1-7}
		Nu. 28 & queen & yes & stairs & 2 & -- & -- \\
		\cmidrule(r){1-7}
		Nu. 38 & queen & yes & stairs & 2 & -- & -- \\
		\cmidrule(r){1-7}
		Nu. 39 & queen & yes & stairs & 2 & yes & -- \\
		\cmidrule(r){1-7}
		Nu. 54 & minor wife & yes & stairs & 1 & -- & -- \\
		\cmidrule(r){1-7}
		Nu. 55 & queen & yes & stairs & 2 & -- & -- \\
		\cmidrule(r){1-7}
		Nu. 57 & queen & yes & stairs & 2 & yes & -- \\
		\cmidrule(r){1-7}
		Nu. 5 & king & yes & stairs & 3 & -- & coffin bench \\
		\cmidrule(r){1-7}
		Nu. 26 & queen & yes & stairs & 2 & yes & -- \\
		\cmidrule(r){1-7}
		Nu. 45 & queen & -- & stairs & 1 & yes & -- \\
		\cmidrule(r){1-7}
		Nu. 18 & king & yes & stairs & 3 & -- & -- \\
		\cmidrule(r){1-7}
		Nu. 10 & king & yes & stairs & 3 & yes & coffin bench \\
		\cmidrule(r){1-7}
		Nu. 25 & queen & yes & stairs & 2 & -- & -- \\
		\cmidrule(r){1-7}
		Nu. 7 & king & yes & stairs & 3 & -- & coffin bench \\
		\cmidrule(r){1-7}
		Nu. 30 & queen & yes & stairs & 2 & -- & coffin bench \\
		\cmidrule(r){1-7}
		Nu. 2 & king & yes & stairs & 3 & yes & coffin bench \\
		\cmidrule(r){1-7}
		Nu. 47 & minor wife & yes & stairs & 2 & yes & coffin bench \\
		\cmidrule(r){1-7}
		Nu. 50 & minor wife & yes & stairs & 1 & -- & -- \\
		\cmidrule(r){1-7}
		Nu. 4 & king & yes & stairs & 3 & yes & coffin bench \\
		\cmidrule(r){1-7}
		Nu. 29 & queen & yes & stairs & 2 & -- & -- \\
		\cmidrule(r){1-7}
		Nu. 52 & minor wife & yes & stairs & 1 & -- & -- \\
		\cmidrule(r){1-7}
		Nu. 19 & king & yes & stairs & 3 & -- & coffin bench \\
		\cmidrule(r){1-7}
		Nu. 46 & minor wife & yes & stairs & 2 & -- & coffin bench \\
		\cmidrule(r){1-7}
		Nu. 49 & minor wife & yes & stairs & 1 & yes & coffin bench \\
		\cmidrule(r){1-7}
		Nu. 11 & king & yes & stairs & 3 & yes & coffin bench \\
		\cmidrule(r){1-7}
		Nu. 31 & queen & yes & stairs & 2 & -- & coffin bench \\
		\cmidrule(r){1-7}
		Nu. 32 & queen & yes & stairs & 2 & -- & coffin bench \\
		\cmidrule(r){1-7}
		Nu. 16 & king & yes & stairs & 3 & -- & coffin bench \\
		\cmidrule(r){1-7}
		Nu. 12 & king & yes & stairs & 3 & yes & -- \\
		\cmidrule(r){1-7}
		Nu. 33 & queen & yes & stairs & 2 & -- & coffin bench \\
	\end{xltabular}
}


\subsection{Object types}


\begin{figure}[H]
	\centering
	\includegraphics[width=1\textwidth, height=12cm, keepaspectratio]{MN-royal-objs}
	\caption{Distribution of object types in Middle Napatan royal tombs by occupant's title.}
	\label{fig:MN-royal-objs}
\end{figure}


Shabtis remained the most common object type, continuing the trend established in the 25th Dynasty phase (Figure \ref{fig:MN-royal-objs}). The distribution, however, is different; with 1243 shabtis found in queens' tombs, 641 in kings' and 8 in minor wives'. Ceramic vessels are the second most prevalent type, comprising 353 associated with queens, 280 with kings and 88 with minor wives. A significant shift occurred in the diversity of object types. Unlike earlier phases, when queens' assemblages were the most heterogeneity, kings' tombs now display the greatest diversity. Nevertheless, even the assemblages of minor wives' comprised eleven different object types, albeit in low numbers. They also contained less objects made specifically for the burial, such as heart scarabs, mummy trappings and stelae.
The distribution of objects does not appear to follow a pattern regarding the owner, as one weapon was found with queen. Heart scarabs were found in three tombs. Kings Amani-nataki-lebte and Sikhespiqo (Nu. 10 and 4) were each accompanied by one instance of this object, while queen Atmataka (Nu. 55), wife of the first king of this phase, Aramatelqo, contained two heart scarabs, similar to queen Yeturow from the previous phase.
Additionally, some object types were totally absent despite their occurrence in other phases, namely scarabs and animal parts. In contrast, bead nets were exclusively to queens, although this may be the result of bad preservation of the stringing.


Cylinder sheaths continue to be featured in royal burials, particularly in those of kings, but in reduced quantities. Five kings were buried with this objects and the majority contained one instance. The only king found with two cylinder sheaths was Amaniastabarqo (Nu. 2). Similar to the Early Napatan phase, only one queen had a cylinder sheath, queen of Malonaqen (Nu. 5)


\subsection{Amulets}


\subsubsection{Amulet distribution}


A significant shift in the distribution of amulets is evident during the Middle Napatan period (Table \ref{tab:MN-royal-distro}). Notably, amulets are no longer present in all king's burials. They were included in only 30\% of tombs. On the other hand, 72\% of queen's graves included at least one amulet. Out of the six burials of minor wives, one contained a single amulet. 
Furthermore, the median number of amulets for kings is much lower than the previous phases, only one. Those from queens also decreased to 0, while the median for minor wives remained at 0.


{\small
	\begin{xltabular}{\textwidth}{X X X X X X}
		\caption{Distribution of amulets in Middle Napatan royal tombs by owner's title, including total counts of amulets and tombs, tombs containing amulets and amulet frequency (\%) and median.}
		\label{tab:MN-royal-distro} \\
		\toprule
		\textbf{Owner} & \textbf{Total number of tombs} & \textbf{Total number of amulets} & \textbf{N. of tombs with amulets} & \textbf{\% of tombs with amulets} & \textbf{Median n. of amulets per tomb} \\
		\midrule
		\endfirsthead
		\multicolumn{6}{c}{\textit{Continued from previous page}} \\
		\toprule
		\textbf{Owner} & \textbf{Total number of tombs} & \textbf{Total number of amulets} & \textbf{N. of tombs with amulets} & \textbf{\% of tombs with amulets} & \textbf{Median n. of amulets per tomb} \\
		\midrule
		\endhead
		\bottomrule
		\endlastfoot
		King & 11 & 37 & 6 & 30\% & 1 \\
		\cmidrule(r){1-6}
		Queen & 13 & 18 & 4 & 72\% & 0 \\
		\cmidrule(r){1-6}
		Minor wife & 6 & 1 & 1 & 17\% & 0 \\
		\midrule
		\textbf{Total} & \textbf{30} & \textbf{56} & \textbf{11} & -- & -- \\
	\end{xltabular}
}


\subsubsection{Amulet types}


\begin{figure}[H]
	\centering
	\includegraphics[width=1\textwidth]{MN-royal-types}
	\caption{Distribution of amulet types in Middle Napatan royal tombs by occupant's title.}
	\label{fig:MN-royal-types}
\end{figure}


The diversity of amulet types also decreases compared to the phases, with a total absence of object and miscellaneous types of amulets (Figure \ref{fig:MN-royal-types}). Symbols continue to be the most popular category (39), primarily associated with kings (29) and queens (10). Less common, with only 12 instances, are animal amulets. They are almost equally distributed between queens (6) and kings (5) and constitute the only amulet type present in the burial of a minor wife. Notably, only a single deity amulet was attested, underscoring a continued decline in their use that started with the 25th Dynasty tombs at Nuri. Nevertheless, queens' burials continue to display the most diverse set of amulet types, with all four categories.


\subsubsection{Amulet motifs}


\begin{figure}[H]
	\centering
	\includegraphics[width=1\textwidth]{MN-royal-motifs}
	\caption{Distribution of amulet motifs in Middle Napatan royal tombs by occupant's title.}
	\label{fig:MN-royal-motifs}
\end{figure}


\begin{wrapfigure}{l}{0.3\textwidth}
	\centering
	\includegraphics[width=1\linewidth, height=5cm, keepaspectratio, clip, trim={10cm 10cm 10cm 10cm}]{MN-20.290}
	\caption{Gold amulet depicting a \textit{ba} bird (MFA 20.290).}
	\label{fig:MN-ba}
\end{wrapfigure} 


Queens continued to showcase the most diverse assemblage of amulet motifs, with twelve different ones, against seven from kings (Figure \ref{fig:MN-royal-motifs}). \textit{Menat} amulets remained the most popular motif (30 examples in total), and were mostly found in kings' tombs (28), and two instances with queens. The corpus includes an amulet representing the god Hapi, one of the Four Son of Horus. However, amulets of the other three gods are missing. It is possible, as we saw with the Early Napatan amulets, that this amulet was part of a set that was stolen. Notably, the tomb of king Sichaspiqo is the only one with an amulet representing a \textit{ba} bird with outspread wings made of gold (Figure \ref{fig:MN-ba}).

 
Among the animal amulets are vultures, snakes, rams, one cow and one hawk. The ram amulets display design variations (Figure \ref{fig:MN-rams}). One of them is a ram's head with sun disk and double \textit{uraeus}. It is notable that this type of amulet is used only during the Middle Napatan period, as depictions of the double \textit{uraeus} in Kushite royal crowns are prevalent during the 25th Dynasty period \parencites{labiaGoddessQueenGoddess2023, coxImperialArtDuality2024}. Another amulet depicts a ram with four heads—a form associated with the four-headed ram god Banebdjedet (Figure \ref{fig:MN-rams}, right). This four-headed ram amulet is the only one found across all royal burials, associated with a minor wife of Aramatelqo. Contemporary amulets with a similar design are known from Egypt \parencites[fig.~41f]{andrewsAmuletsAncientEgypt1994}.
 
 
 \begin{figure}[H]
 	\centering
 	\begin{minipage}{0.3\textwidth}
 		\includegraphics[width=\textwidth, height=4cm, keepaspectratio]{MN-20.268}
 	\end{minipage}
 	\hspace{-25mm}%
 	\begin{minipage}{0.3\textwidth}
 		\includegraphics[width=\textwidth, height=4cm, keepaspectratio]{MN-20.269}
 	\end{minipage}
 	\hspace{-25mm}%
 	\begin{minipage}{0.35\textwidth}
 		\includegraphics[width=\textwidth, height=4cm, keepaspectratio, clip, trim={8cm 6cm 8cm 9cm}]{MN-24.1017}
 	\end{minipage}
 	\caption{Middle Napatan ram amulets (MFA  20.268, 20.269 and 24.1017).}
 	\label{fig:MN-rams}
 \end{figure}
 
 
 Amulets stamped onto gold sheets continue to be produced during the Middle Napatan period (Figure \ref{fig:MN-pics}, \ref{fig:MN-pics2}). These amulets could represent a variety of amulet types, such as animals, symbols and nature elements. Among the animals are a cow wearing the Hathor headdress, a snake and vultures. One of the vulture amulets stands out due to its depiction of two vultures side-by-side with a single suspension ring on the back (Figure \ref{fig:MN-pics}, bottom). This could suggest that the amulet was left unfinished and uncut at the time it was included as a grave good. The remaining symbols comprise the Isis knot, the hieroglyphic signs \textit{akhet} and \textit{nefer}, \textit{uraeus}, \textit{ba} bird and, finally, two \textit{udjat} amulets, a sharp decline from previous phases where it had been one of the most popular motifs. Nature elements depict two papyrus columns, a lotus flower and a double leaf (Figure \ref{fig:MN-pics2}).
 
 
 \begin{figure}[H]
 	\centering
 	\begin{minipage}{0.2\textwidth}
 		\includegraphics[width=\linewidth, height=4cm, keepaspectratio]{MN-20.347}
 	\end{minipage}
 	\hspace{-2mm}%
 	\begin{minipage}{0.248\textwidth}
 		\includegraphics[width=\linewidth, height=4cm, keepaspectratio, clip, trim={5cm 0cm 5cm 0cm}]{MN-20.351}
 	\end{minipage}
 	\vspace{-5mm}%
 	\begin{minipage}{0.3\textwidth}
 		\includegraphics[width=\linewidth, height=4cm, keepaspectratio]{MN-20.348}
 	\end{minipage}
 	\hspace{-2mm}%
 	\begin{minipage}{0.47\textwidth}
 		\includegraphics[width=\linewidth, height=4cm, keepaspectratio, clip, trim={3cm 5cm 3cm 3cm}]{MN-20.1061}
 	\end{minipage}
 	\caption{Middle Napatan amulets depicting animals: cow (MFA 20.347), snake (MFA 20.351), and vultures (MFA 20.348 and 20.1061).}
 	\label{fig:MN-pics}
 \end{figure}
 
 
 \begin{figure}[H]
 	\centering
 	\begin{minipage}{0.17\textwidth}
 		\includegraphics[width=\textwidth, height=4cm, keepaspectratio]{MN-20.270}
 	\end{minipage}
 	\hspace{-5mm}%
 	\begin{minipage}{0.18\textwidth}
 		\includegraphics[width=\textwidth, height=4cm, keepaspectratio]{MN-20.346}
 	\end{minipage}
 	\hspace{-1mm}%
 	\begin{minipage}{0.18\textwidth}
 		\includegraphics[width=\textwidth, height=4cm, keepaspectratio]{MN-20.352}
 	\end{minipage}
 	\hspace{-8mm}%
 	\begin{minipage}{0.22\textwidth}
 		\includegraphics[width=\textwidth, height=4cm, keepaspectratio]{MN-20.366}
 	\end{minipage}
 	\hspace{-5mm}%
 	\begin{minipage}{0.18\textwidth}
 		\centering
 		\includegraphics[width=\textwidth, height=4cm, keepaspectratio]{MN-20.369}
 	\end{minipage}%
 	\caption{Gold amulets representing a double leaf (MFA 20.270), \textit{uraeus} (MFA 20.346), \textit{nefer} (MFA 20.352), lotus flower (MFA 20.366) and \textit{tjet} (MFA 20.369).}
 	\label{fig:MN-pics2}
 \end{figure}
 
 
 \subsubsection{Amulet materials}
 
 
 \begin{figure}[H]
 	\centering
 	\includegraphics[width=1\textwidth]{MN-royal-materials}
 	\caption{Distribution of amulet materials in Middle Napatan royal tombs by occupant's title.}
 	\label{fig:MN-royal-materials}
 \end{figure}
 
 
 The range of materials used for the manufacture of amulets also decreased, with five or six total materials (Figure \ref{fig:MN-royal-materials}). This number depends on whether beryl is classified as amazonite, as discussed earlier. If true, amazonite accounts for a total of 3 amulets, two from kings' graves and one from a queen's. Faience remained the most prevalent material, found in all types of tombs, the majority with kings (28), then with queens (2) and minor wives (1). It was followed by gold, with 6 instances with kings and 12 with queens. Jasper was the composition of one amulet of a king and two of queens, while the sole travertine amulet was found with a queen.
 
 Imported materials include jasper and travertine, both from the Egyptian Eastern Desert, and amazonite, possibly from Ethiopia. The scarcity of further material composition, specially those sourced in the Eastern Desert, could possibly indicate that, starting from the Early Napatan phase, access to trade routes were more limited when compared to the 25th Dynasty.
 
 
 \section{Late Napatan}
 
 
 \subsection{Tomb structure}
 
 
 \begin{figure}[H]
 	\centering
 	\includegraphics[width=1\textwidth]{LN-royal-tombs}
 	\caption{Distribution of superstructure and substructure types by tomb occupant's title in Late Napatan royal burials.}
 	\label{fig:LN-royal-tombs}
 \end{figure}
 
 
Late Napatan tombs were constructed both in El-Kurru and Nuri, though, pyramid Ku. 1, intended for a king's burial, was never used for its purpose (Figure \ref{fig:LN-royal-tombs}). The adjacent pyramid Ku. 2 is linked to his queen and both tombs were built when the royal cemetery had already moved to Nuri \parencites{emberlingExcavationPyramidKu2015}. Their inclusion in the analysis is justified by their completed structure, despite the absence of a burial. Unlike earlier phases, all superstructures in this period are identifiable as pyramids, including those of minor wives. This is either a result of better preservation conditions or a progression of superstructures to be only pyramids. It is important to note that, even though we know very little about this specific phase during the Napatan period, the predominance of pyramids could indicate that the kingdom was prospering and could acquire the human and material resources for all royalty to be buried in pyramids. Similarly, all substructures consist of chambers.
 
Similarly to Middle Napatan, the majority of tombs featured a chapel, with the sole exception of Ku. 2, and stairs continued to be the only form of entrance into the burial (Table \ref{tab:LN-royal-tomb-elements}). The number of chambers follows a clear pattern: all kings' pyramids feature three chambers, while queens' have two. The two sole exceptions are the tombs of Baskakeren's (Nu. 17)---with two chambers---, and, possibly, Akhratan's (Nu. 14), as it was not excavated due to dangerous conditions. Minor wives' tombs could had more flexibility, containing either one or two chambers. While all kings' tombs bore decoration (either in chapels or chambers), decoration was found in only two queens' tombs and was entirely absent from minor wives' tombs. The majority of tombs included a coffin bench (10 out of 14), however not all kings' burials and queens' had one, whereas all of the minor wives' tombs did. All tombs featuring eye inlays for a coffin include a bench as well.
 
 
{\small
 	\begin{xltabular}{\textwidth}{X X X X X X X}
 		\caption{Distribution of architectural elements by tomb owner's title in Late Napatan royal burials.}
 		\label{tab:LN-royal-tomb-elements} \\
 		\toprule
 		\textbf{Tomb} & \textbf{Owner} & \textbf{Chapel} & \textbf{Type of entrance} & \textbf{N. of chambers} & \textbf{Decoration} & \textbf{Bed holes and/or coffin bench} \\
 		\midrule
 		\endfirsthead
 		\multicolumn{7}{c}{\textit{Continued from previous page}} \\
 		\toprule
 		\textbf{Tomb} & \textbf{Owner} & \textbf{Chapel} & \textbf{Type of entrance} & \textbf{N. of chambers} & \textbf{Decoration} & \textbf{Bed holes and/or coffin bench} \\
 		\midrule
 		\endhead
 		\bottomrule
 		\endlastfoot
 		Nu. 17 & king & yes & stairs & 2 & yes & coffin bench \\
 		\cmidrule(r){1-7}
 		Nu. 13 & king & yes & stairs & 3 & yes & -- \\
 		\cmidrule(r){1-7}
 		Nu. 34 & queen & yes & stairs & 2 & yes & coffin bench \\
 		\cmidrule(r){1-7}
 		Nu. 44 & queen & yes & stairs & 2 & -- & coffin bench \\
 		\cmidrule(r){1-7}
 		Nu. 61 & queen & yes & stairs & 2 & -- & coffin bench \\
 		\cmidrule(r){1-7}
 		Ku. 1 (unused) & king & yes & stairs & 3 & yes & coffin bench \\
 		\cmidrule(r){1-7}
 		Ku. 2 & queen & -- & stairs & 2 & -- & -- \\
 		\cmidrule(r){1-7}
 		Nu. 43 & minor wife & yes & stairs & 2 & -- & coffin bench \\
 		\cmidrule(r){1-7}
 		Nu. 48 & minor wife & yes & stairs & 2 & -- & coffin bench \\
 		\cmidrule(r){1-7}
 		Nu. 51 & minor wife & yes & stairs & 1 & -- & coffin bench \\
 		\cmidrule(r){1-7}
 		Nu. 14 & king & yes & stairs & not excavated & yes & -- \\
 		\cmidrule(r){1-7}
 		Nu. 37 & queen & yes & stairs & 2 & -- & coffin bench \\
 		\cmidrule(r){1-7}
 		Nu. 15 & king & yes & stairs & 3 & yes & -- \\
 		\cmidrule(r){1-7}
 		Nu. 56 & queen & yes & stairs & 2 & yes & coffin bench \\
 	\end{xltabular}
}
 
 
 \subsection{Object types}
 
 
 \begin{figure}[H]
 	\raggedleft
 	\includegraphics[width=0.4\textwidth, height=10cm, keepaspectratio]{LN-royal-objs-el-kurru}
 	\raggedright
 	\includegraphics[width=0.59\textwidth, height=10cm, keepaspectratio]{LN-royal-objs-nuri}
 	\caption{Distribution of object types in Late Napatan royal tombs by occupant's title.}
 	\label{fig:LN-royal-objs}
 \end{figure}
 
 
 Object types at the two burials in El-Kurru were very limited (Figure \ref{fig:LN-royal-objs}). The unused king's tomb (Ku. 1) only yielded ceramic and offering vessels, whereas the queen's pyramid (Ku. 2) featured 152 shabtis\footnote{Reisner found the shabtis in Ku. 6 and presumed that they were originally from Ku. 2 \parencites[p.~25, 43]{dunhamElKurru1950}}. At Nuri, shabtis are much less numerous than previous periods, with 34 instances associated with queens and 25 with kings. Scarabs and animals parts, absent during the Middle Napatan phase, resumed in Late Napatan burial assemblages. In contrast, furniture,  scarabs and statuettes are entirely absent in these later tombs. The disappearance of furniture is likely attributable to poorer preservation conditions, given that such items were often made of perishable materials, such as wood. The case of heart scarabs is more intriguing; although always a rare occurrence in previous phases, they were nonetheless present, sometimes in a pair. Their complete absence in the Late Napatan period is therefore significant and may reflect a shift in burial practices.
 
 
 A unique find from this phase is a model Osiris whip (or flail) made of gold, found in the tomb of king Harsijotef (Nu. 13). As no other examples appear in any previous phase, its singular occurrence may be a result of looting if others were also made of gold. Flails have been found in the tombs of Egyptian kings, namely Senebtisi and Hapiankhtifi from the Middle Kingdom \parencites[p.~78, figs.~35, 9]{maceTombSenebtisiLisht1916}[p.~244, cat.~181]{quirkeUnderstandingDeathJourney2015}. Tutmosis III's mummy also suggested that he was buried grasping a cylindrical object in his left hand, likely the flail, which was taken when his burial was looted \parencites[p.~5]{smithRoyalMummiesCG1912}. Bead nets are an interesting case as well; since their emergence in the 25th Dynasty, they are exclusively found in female burials, primarily at queens'. This pattern suggests that they may have a gender association.
 
 
 \subsection{Amulets} 
 
 
 \subsubsection{Amulet distribution}
 
 
 {\small
 	\begin{xltabular}{\textwidth}{X X X X X X}
 		\caption{Distribution of amulets in Late Napatan royal tombs by owner's title, including total counts of amulets and tombs, tombs containing amulets and amulet frequency (\%) and median.}
 		\label{tab:LN-royal-distro} \\
 		\toprule
 		\textbf{Owner} & \textbf{Total number of tombs} & \textbf{Total number of amulets} & \textbf{N. of tombs with amulets} & \textbf{\% of tombs with amulets} & \textbf{Median n. of amulets per tomb} \\
 		\midrule
 		\endfirsthead
 		\multicolumn{6}{c}{\textit{Continued from previous page}} \\
 		\toprule
 		\textbf{Owner} & \textbf{Total number of tombs} & \textbf{Total number of amulets} & \textbf{N. of tombs with amulets} & \textbf{\% of tombs with amulets} & \textbf{Median n. of amulets per tomb} \\
 		\midrule
 		\endhead
 		\bottomrule
 		\endlastfoot
 		King & 5 & 8 & 3 & 60\% & 1 \\
 		\cmidrule(r){1-6}
 		Queen & 6 & 9 & 3 & 50\% & 0 \\
 		\cmidrule(r){1-6}
 		Minor wife & 3 & 0 & 0 & 0\% & 0 \\
 		\midrule
 		\textbf{Total} & \textbf{14} & \textbf{17} & \textbf{6} & -- & -- \\
 	\end{xltabular}
 }
 
 
 The Late Napatan period featured the least amount of amulets of all Napatan phases, with only 17 amulets distributed in 6 tombs (Table \ref{tab:LN-royal-distro}). As in the Middle Napatan period, not all kings' tombs contained amulets. However, their frequency increases, with amulets appearing in 60\% of kings' burials. 50\% of queens' were buried with amulets, in contrast to the total absence of amulets in the burials of minor wives. Furthermore, queens and kings were buried with the same average number of amulets.
 
 
 \subsubsection{Amulet types}
 
 \begin{figure}[H]
 	\centering
 	\includegraphics[width=1\textwidth]{LN-royal-types}
 	\caption{Distribution of amulet types in Late Napatan royal tombs by occupant's title.}
 	\label{fig:LN-royal-types}
 \end{figure}
 
 
 There are only two categories of Late Napatan amulets: symbols and deities (Figure \ref{fig:LN-royal-types}). Symbols continue to be more numerous, with 15 amulets: seven with kings and eight with queens. Deities comprise two amulets, one in each type of burial.
 
 
 \subsubsection{Amulet motifs}
 
 \begin{figure}[H]
 	\centering
 	\includegraphics[width=1\textwidth]{LN-royal-motifs}
 	\caption{Distribution of amulet motifs in Late Napatan royal tombs by occupant's title.}
 	\label{fig:LN-royal-motifs}
 \end{figure}
 
 
 Unlike previous phases, in Late Napatan, kings' burials feature the most diverse set of amulet motifs, with six different ones, in contrast with only two for queens (Figure \ref{fig:LN-royal-motifs}). The \textit{udjat} resumed its position as the most common amulet, primarily associated with queens. Both of the kings' \textit{udjat} amulets were buried with Harsijotef (Nu. 13) and were strung in a gold wire forming bracelets (Figure \ref{fig:LN-royal-udjat}). Conversely, the \textit{menat} is represented by only one amulet, included in the tomb of king Baskakeren (Nu. 17).
 
 
 \begin{figure}[H]
 	\centering
 	\begin{minipage}{0.49\textwidth}
 		\centering
 		\includegraphics[width=\textwidth, height=5cm, keepaspectratio]{LN-20.324}
 	\end{minipage}
 	\hspace{-40mm}%
 	\begin{minipage}{0.49\textwidth}
 		\centering
 		\includegraphics[width=\textwidth, height=5cm, keepaspectratio]{LN-20.325}
 	\end{minipage}
 	\caption{\textit{Udjat} amulets in gold wire bracelets (MFA 20.324, 20.325).}
 	\label{fig:LN-royal-udjat}
 \end{figure}
 
 
 The \textit{uraeus}, previously exclusive to the tombs of queens and minor wives, now appeared solely in the burials of kings. In contrast, the association of the Four Sons of Horus with queens continued, exemplified by a single amulet of Qebehsenuef found in a queen's burial—although a complete set of all four deities remained absent. The second-known Anubis amulet from a royal burial was found in the tomb of king Nastasen. Both this one and the first (from a 25th Dynasty minor wife at Nuri) were made of gold. Similarly, \textit{ba} bird amulets remained very rare, with only three instances recorded across the entire royal burial corpus; each was found in a different burial (both king and queen) from a different phase (25th Dynasty, Middle and Late Napatan).
 
 
 \begin{wrapfigure}{l}{0.3\textwidth}
 	\centering
 	\includegraphics[width=1\linewidth, height=5cm, keepaspectratio]{LN-griffin}
 	\caption{Gold amulet depicting a winged quadruple animal, possibly a griffin \parencites[pl.~CXX.A, 17-2-1978]{dunhamNuri1955}}
 	\label{fig:LN-griffin}
 \end{wrapfigure}
 
 
 The winged griffin is a particularly interesting motif (Figure \ref{fig:LN-griffin}). Unfortunately, I was not able to locate a recent picture of it, even though it was assigned to the Sudan National Museum. In the report, Dunham included a question mark after the description, likely because he was not sure of the design. Examining the photograph included in the report, it appears to depict a winged quadruped animal with an avian head, which could be interpreted as a griffin. A similar motif was found in Meroe South and will be discussed in Chapter 5. The iconography of the griffin is attested in Egypt with many different designs, including winged \parencites{hsuGriffinVisualWritten2011}. Depictions of winged griffins are attested in the Middle Kingdom, with one interpretation pointing their connection to magical knives and the rebirth of the sun god \parencites{colloffReinterpretingGriffinsMiddle2022}, while others suggest their connections with the god Nefertum \parencites{sabbahyMiddleBronzeAge2017}. From the New Kingdom, the ceremonial axe of Ahmose contains a winged griffin above a scene of the king striking an enemy, prompting the interpretation that the griffin is a representation of the king \parencites{wyattGraspingGriffinIdentifying2010}. However, the occurrences of this motif as an amulet appear to be unique to Nubia during this period.
 
 \subsubsection{Amulet materials}
 
 \begin{figure}[H]
 	\centering
 	\includegraphics[width=1\textwidth]{LN-royal-materials}
 	\caption{Distribution of amulet materials in Late Napatan royal tombs by occupant's title.}
 	\label{fig:LN-royal-materials}
 \end{figure}
 
 The trend towards reduced material diversity culminated in this phase, with amulets being made of only four distinct materials (Figure \ref{fig:LN-royal-materials}). Most strikingly, faience was not the most popular material, with only a single \textit{menat} example. The \textit{udjat} amulets were crafted from serpentinite and milky quartz, while gold became the exclusive material for all other motifs and the most prevalent material. The two \textit{udjat} amulets were made of serpentinite and milky quartz, which were sourced from the Aswan region or further north \parencites[p.~506-507, 759]{harrellArchaeologyGeologyAncient2024}. This marked decline in amulet use and diversity may reflect the broader decline of the Napatan period.
 
 \section{Conclusion}
 
 The primary aim of this chapter was to employ chronological analysis and comparison between each phase of the Napatan period in relation to tomb structure, funerary objects and amulets. A quantitative analysis of the royal burial sites shows a clear evolution of the tomb configuration. Regarding super and substructures, tumuli is succeeded by pyramids and the pit by underground chambers. Due to the inclusion of chambers, stairs also become ubiquitous. Middle and Late Napatan tombs have more chapels, with some exceptions. Although destruction of earlier chapels must be considered, it is notable that they survived more in later burials. From the Early Napatan phase, there are no more cutouts for burial beds and coffin benches occur in the majority of tombs. Decoration is sparse and does not seem to follow a gender or title association. It is likely that the lack of decoration is a result of poor preservation. A clear hierarchy based on gender and title in evident in the number of chambers.  Three-chambered tombs, introduced for kings in the Early Napatan phase, became the standard in the Middle Napatan phase. In contrast, two chambers are associated with queens and one chamber with minor wives, though exceptions exist in both cases.
 
 
 Funerary object assemblages show continuity despite variation across phases and titles. Only five object types were exclusive to a single period: pebbles, games and toys, shrines, and likely poorly-preserved clothing from the 25th Dynasty; as well as the Osiris whip from Late Napata. In contrast, seven types occurred in all phases, and eight objects introduced in the 25th Dynasty became permanent additions. These objects include bead nets, burial containers, canopics, mummy trappings, staffs, stelae, toiletries and tools. This suggests the Kushite rule of Egypt prompted a change in the funerary cult, as evidenced by the 25th Dynasty's diverse assemblage. Crucially, the Kushites were not passive recipients of an Egyptian tradition but active innovators, introducing new objects like bead nets and re-introducing older practices such as offering tables. On the other hand, cylinder sheaths represent a uniquely Napatan innovation. First adopted in the Early Napatan phase following the withdrawal from Egypt, they persisted as a distinctive feature of royal burials in later tombs.
 
 
 The 25th Dynasty represents the high point of Napatan amulet use, featuring the greatest diversity of motifs, the highest quantity, and unique designs that blend Egyptian and local traditions. Following this peak, both the number and variety of amulets declined significantly. Symbols are the most numerous category of all amulets, on account of \textit{udjat} and \textit{menat} amulets, the most popular motifs across all phases. In fact, the \textit{udjat} is the only amulet motif occurring in all phases. This motif was extremely popular in Egypt since the Old Kingdom and was exported to other locations such as the Levant and Israel \parencites[p.~43-44]{andrewsAmuletsAncientEgypt1994}[p.~611-773]{herrmannÄgyptischeAmuletteAus1994}{arroyoCulturaMaterialAntigo2024}. The material composition of amulets indicate a strong preference for faience, possibly due to the availability of the material. Yet, the diversity of materials also suggests access, albeit restricted in later phases, to a range of other materials, including imported ones, mostly from the Egyptian Eastern Desert region.
 